%% toc: P3 35 35.5,35.6,35.7
\section{Proofs: infinitary languages, Karp's theorem, Scott's theorem and model existence}
\label{sec:annex-p03-w03-p03-039-marker}

Structures are nonempty. For an assignment \(s\) and an element \(e\),
\(s[w\mapsto e]\) agrees with \(s\) except that it sends \(w\) to \(e\). The
truth of a formula under an assignment depends only on the values of its
free variables; this is checked by the same induction on formulas as in
first-order logic, the new cases \(\bigwedge,\bigvee\) being immediate.

\begin{annexitem}{annex:p03-w03-p03-039-iso-invariance}{thm:p03-w03-p03-039-iso-invariance}{Isomorphism preserves infinitary truth}
\begin{proof}
By induction on \(\varphi\). For a term \(t\), \(h(t^{\mathcal M}[s])=t^{\mathcal N}[h\circ s]\)
by induction on \(t\), because \(h\) commutes with the interpretations of
the function and constant symbols; with the fact that \(h\) preserves and
reflects the relations and equality (it is a bijection) this gives the
atomic case. The cases \(\neg\), \(\bigwedge_{\varphi\in X}\) and
\(\bigvee_{\varphi\in X}\) follow from the induction hypothesis for the
members of \(X\). If \(\mathcal M,s\models\exists w\,\psi\), choose \(e\in M\)
with \(\mathcal M,s[w\mapsto e]\models\psi\); by the induction hypothesis
\(\mathcal N,h\circ s[w\mapsto h(e)]\models\psi\), and
\(h\circ(s[w\mapsto e])=(h\circ s)[w\mapsto h(e)]\). Conversely a witness
\(e'\in N\) for \(\mathcal N,h\circ s\models\exists w\,\psi\) has the form
\(h(e)\) because \(h\) is surjective, and the same identity applies. The
case \(\forall w\) is the same argument with ``every'' in place of ``some'',
again using surjectivity. For a sentence the assignment is irrelevant.
\end{proof}
\end{annexitem}

\begin{annexitem}{annex:p03-w03-p03-039-compactness-failure}{thm:p03-w03-p03-039-compactness-failure}{Compactness fails for \(\mathcal L_{\omega_1,\omega}\)}
\begin{proof}
Let \(\Delta\subseteq\Gamma\) be finite. Then only finitely many sentences
\(d\ne c_i\) lie in \(\Delta\); choose \(j\) with \(d\ne c_j\notin\Delta\).
Interpret the structure on \(\omega\) by \(c_i\mapsto i\) and \(d\mapsto j\).
Every element \(m\) equals the interpretation of \(c_m\), so the sentence
\(\forall v\bigvee_iv=c_i\) holds, and \(d\ne c_i\) holds for every
\(i\ne j\), in particular for every such sentence in \(\Delta\). So
\(\Delta\) has a model. If \(\mathcal M\models\Gamma\), then
\(\mathcal M\models\bigvee_iv=c_i\) at the element \(d^{\mathcal M}\), so
\(d^{\mathcal M}=c_i^{\mathcal M}\) for some \(i\), contradicting
\(\mathcal M\models d\ne c_i\). The set \(\Gamma\) consists of
\(\mathcal L_{\omega_1,\omega}\)-sentences, hence of
\(\mathcal L_{\infty,\omega}\)-sentences.
\end{proof}
\end{annexitem}

\begin{annexitem}{annex:p03-w03-p03-039-formal-negation}{thm:p03-w03-p03-039-formal-negation}{Formal negation is negation}
\begin{proof}
The recursion defining \({\sim}\) follows the formation of formulas, which is
well founded, so the claims may be proved by induction on \(\varphi\). For
atomic \(\varphi\) the formula \({\sim}\varphi\) is \(\neg\varphi\). If
\(\varphi=\neg\theta\), then \({\sim}\varphi=\theta\), and \(\neg\neg\theta\)
holds exactly when \(\theta\) does. If \(\varphi=\bigwedge_{\theta\in X}\theta\),
then \(\mathcal M,s\models\neg\varphi\) iff some \(\theta\in X\) fails iff
some \(\theta\in X\) has \(\mathcal M,s\models\neg\theta\) iff, by the
induction hypothesis, some \(\theta\in X\) has
\(\mathcal M,s\models{\sim}\theta\), iff
\(\mathcal M,s\models\bigvee_{\theta\in X}{\sim}\theta={\sim}\varphi\). The
case \(\bigvee\) is dual. If \(\varphi=\exists w\,\theta\), then
\(\neg\varphi\) holds at \(s\) iff \(\neg\theta\) holds at
\(s[w\mapsto e]\) for every \(e\), iff (induction hypothesis) \({\sim}\theta\)
holds at every \(s[w\mapsto e]\), iff \(\forall w\,{\sim}\theta={\sim}\varphi\)
holds at \(s\); the case \(\forall\) is dual. The same induction shows that
\({\sim}\) does not change free variables, since it never changes an atomic
formula's variables and never adds or removes a quantifier, and that
substituting a term \(t\) for the free occurrences of \(v\) commutes with
\({\sim}\), since both operations act clause by clause and \({\sim}\) leaves
atomic formulas inside a \(\neg\).
\end{proof}
\end{annexitem}

\begin{annexitem}{annex:p03-w03-p03-039-countable-fragment}{thm:p03-w03-p03-039-countable-fragment}{Countable sets of sentences lie in countable fragments}
\begin{proof}
First, every \(\mathcal L_{\omega_1,\omega}\)-formula has countably many
subformulas and countably many variables: by induction, an atomic formula
has finitely many, and the other formation steps take a countable union of
countable sets together with one new formula or variable.

Let \(S_0\) be \(T\) together with all atomic formulas of \(\tau\) whose
variables lie in \(V\). Since \(\tau\) and \(V\) are countable there are
countably many terms over them, hence countably many such atomic formulas,
and \(S_0\) is countable. Given a countable \(S_k\), let \(S_{k+1}\) be
\(S_k\) together with: all subformulas of members of \(S_k\); all results
of substituting a term over \(\tau\) and \(V\) for a free variable of a member
of \(S_k\); all \({\sim}\varphi\) for \(\varphi\in S_k\); and all
\(\neg\varphi\), \(\varphi\land\psi\), \(\varphi\lor\psi\),
\(\exists v\,\varphi\), \(\forall v\,\varphi\) with \(\varphi,\psi\in S_k\),
\(v\in V\). Each of these is a countable union of countable sets, so
\(S_{k+1}\) is countable. Put \(A=\bigcup_kS_k\), which is countable and
consists of \(\mathcal L_{\omega_1,\omega}\)-formulas.

Every variable occurring in \(A\) is in \(V\): this holds in \(S_0\) by the
choice of \(V\), and each operation above only uses variables of the input
formulas, of terms over \(V\), or of \(V\) (\({\sim}\) introduces none, by
\cref{thm:p03-w03-p03-039-formal-negation}). Each closure condition of
\cref{def:p03-w03-p03-039-fragment} concerns one or two members of \(A\),
which lie in some common \(S_k\), and its output lies in \(S_{k+1}\). So
\(A\) is a fragment with witness \(V\), and \(T\subseteq S_0\subseteq A\).
If \(B\) is any fragment with witness \(V\) and \(T\subseteq B\), then
\(S_0\subseteq B\) by clause~1, and \(S_k\subseteq B\) implies
\(S_{k+1}\subseteq B\) by clauses~2--5; hence \(A\subseteq B\).
\end{proof}
\end{annexitem}

\begin{annexitem}{annex:p03-w03-p03-039-downward-ls}{thm:p03-w03-p03-039-downward-ls}{Downward L\"owenheim--Skolem for fragments}
\begin{proof}
Let \(V\) witness that \(A\) is a fragment. Clause~1 of
\cref{def:p03-w03-p03-039-fragment} puts \(v=v\) in \(A\) for each
\(v\in V\), and an atomic formula for each symbol of \(\tau\), so
\(\lvert V\rvert\le\lvert A\rvert\), \(\lvert\tau\rvert\le\lvert A\rvert\),
and \(\lvert A\rvert\) is infinite.

For each \(\psi\in A\) and \(w\in V\), let \(\bar x\) be the increasing list
of the free variables of \(\psi\) other than \(w\); it is finite by
hypothesis. Using the axiom of choice, fix
\(F_{\psi,w}:M^{\lvert\bar x\rvert}\to M\) such that, whenever
\(\mathcal M\models\exists w\,\psi\) at an assignment sending \(\bar x\) to
\(\bar a\), \(\psi\) holds at the assignment sending \(\bar x\) to \(\bar a\)
and \(w\) to \(F_{\psi,w}(\bar a)\); where no witness exists the value is an
arbitrary element. There are at most \(\lvert A\rvert\) such functions.

Fix \(m_0\in M\). Let \(N_0\) be \(X\cup\{m_0\}\) together with the
interpretations of the constant symbols, and let \(N_{k+1}\) be \(N_k\)
together with all values \(F_{\psi,w}(\bar a)\) and \(f^{\mathcal M}(\bar a)\)
for \(f\) a function symbol and \(\bar a\) from \(N_k\). With
\(\lambda=\max(\lvert A\rvert,\lvert X\rvert)\), an infinite cardinal,
\(\lvert N_0\rvert\le\lambda\), and \(\lvert N_k\rvert\le\lambda\) gives
\(\lvert N_{k+1}\rvert\le\lambda+\lvert A\rvert\cdot\lambda=\lambda\). So
\(N=\bigcup_kN_k\) has \(\lvert N\rvert\le\lambda\), contains \(X\), and is
closed under the functions and constants of \(\mathcal M\); let
\(\mathcal N\) be the substructure on \(N\).

We show by induction on \(\varphi\in A\) that for every assignment \(s\)
into \(N\), \(\mathcal N,s\models\varphi\) iff \(\mathcal M,s\models\varphi\);
the constituents of a member of \(A\) are in \(A\) by clause~2. Atomic
formulas: \(\mathcal N\) is a substructure. Negation, \(\bigwedge\) and
\(\bigvee\): induction hypothesis. Let \(\varphi=\exists w\,\psi\). If
\(\mathcal N,s[w\mapsto e]\models\psi\) for some \(e\in N\), the induction
hypothesis gives the same in \(\mathcal M\). If \(\mathcal M,s\models\varphi\),
let \(\bar x\) be as above for \((\psi,w)\); then \(w\in V\) and
\(e=F_{\psi,w}(s(\bar x))\in N\) satisfies
\(\mathcal M,s[w\mapsto e]\models\psi\), so
\(\mathcal N,s[w\mapsto e]\models\psi\) by the induction hypothesis. Let
\(\varphi=\forall w\,\psi\). If \(\mathcal M,s\models\varphi\), the
induction hypothesis gives \(\mathcal N,s[w\mapsto e]\models\psi\) for every
\(e\in N\). If \(\mathcal N,s\models\varphi\) but
\(\mathcal M,s\not\models\varphi\), then
\(\mathcal M,s\models\exists w\,\neg\psi\); as \(\neg\psi\in A\) by
clause~5 and has the same free variables as \(\psi\),
\(e=F_{\neg\psi,w}(s(\bar x))\in N\) has
\(\mathcal M,s[w\mapsto e]\not\models\psi\), and the induction hypothesis
gives \(\mathcal N,s[w\mapsto e]\not\models\psi\), a contradiction. Thus
\(\mathcal N\prec_A\mathcal M\). If \(\tau\) and \(A\) are countable and
\(X=\varnothing\), then \(\lvert N\rvert\le\aleph_0\).
\end{proof}
\end{annexitem}

\begin{annexitem}{annex:p03-w03-p03-039-chang}{thm:p03-w03-p03-039-chang}{Chang's reduction to first-order omitting types}
\begin{proof}
Use the construction of \cref{def:p03-w03-p03-039-chang-construction}. The
fragment \(A\) is countable by \cref{thm:p03-w03-p03-039-countable-fragment}
and its formulas have finitely many free variables, being subformulas or
substitution instances of the sentences of \(T\) or built from such by
finitary steps (each of the operations used to generate \(A\) preserves
finiteness of the free-variable set, and so does a subformula of a sentence).
Hence \(\tau^*\), \(T^*\) and \(\Gamma\) are countable, and each member of
\(T^*\) is a first-order sentence, each \(\gamma_\varphi\) a set of
first-order formulas in the finitely many variables \(\bar x_\varphi\).
The free variables of a conjunct or disjunct \(\theta\in X\) lie among
those of \(\bigwedge X\) or \(\bigvee X\), and those of \(\theta\) lie among
those of \(\exists w\,\theta\) together with \(w\); so each clause applies
\(R_\theta\) to variables available in it.

(1) Let \(\mathcal M\models T^*\) omit every \(\gamma_\varphi\). We show by
induction on \(\varphi\in A\) that \(\mathcal M\models R_\varphi(\bar a)\)
iff \(\mathcal M\models\varphi(\bar a)\) for every tuple \(\bar a\) assigned
to \(\bar x_\varphi\). Atomic \(\varphi\), \(\neg\theta\), \(\exists w\,\theta\)
and \(\forall w\,\theta\): the biconditional axioms and the induction
hypothesis. Let \(\varphi=\bigwedge_{\theta\in X}\theta\). If
\(R_\varphi(\bar a)\), then \(R_\theta\) holds at \(\bar a\) for each
\(\theta\in X\), so each \(\theta\) holds and \(\varphi(\bar a)\) holds. If
\(\varphi(\bar a)\) holds, each \(R_\theta\) holds at \(\bar a\); were
\(R_\varphi(\bar a)\) false, \(\bar a\) would realize \(\gamma_\varphi\),
which is omitted; so \(R_\varphi(\bar a)\). The case of \(\bigvee\) is dual,
using \(\gamma_\varphi=\{R_\varphi\}\cup\{\neg R_\theta:\theta\in X\}\). For a
sentence \(\varphi\in T\), \(R_\varphi\in T^*\) gives
\(\mathcal M\models\varphi\); since \(\varphi\) is a \(\tau\)-sentence, the
\(\tau\)-reduct satisfies it.

(2) Let \(\mathcal M\models T\) and interpret each \(R_\varphi\) as the set
of tuples satisfying \(\varphi\). The axioms of the first five clauses then
restate satisfaction clauses, and the sixth holds because
\(\mathcal M\models T\). A tuple realizing \(\gamma_\varphi\) for
\(\varphi=\bigwedge X\) would satisfy every \(\theta\in X\) but not
\(\varphi\), which is impossible; dually for \(\bigvee\). So the expansion
models \(T^*\) and omits \(\Gamma\).
\end{proof}
\end{annexitem}

\begin{annexitem}{annex:p03-w03-p03-039-countable-back-forth}{thm:p03-w03-p03-039-countable-back-forth}{Back and forth between countable structures}
\begin{proof}
Let \(P\) be a partial isomorphism system and list
\(M=\{m_0,m_1,\ldots\}\), \(N=\{n_0,n_1,\ldots\}\), with repetitions if a
universe is finite. Choose \(f_0\in P\). Given \(f_{2k}\in P\), clause~1 gives
\(f_{2k+1}\in P\) extending it with \(m_k\in\operatorname{dom}(f_{2k+1})\);
clause~2 gives \(f_{2k+2}\in P\) extending that with
\(n_k\in\operatorname{img}(f_{2k+2})\). The union \(h=\bigcup_jf_j\) is a
function, being the union of a chain of functions, with domain \(M\) and
image \(N\). For an atomic \(\varphi(\bar v)\) and \(\bar a\) from \(M\), all
of \(\bar a\) lies in some \(\operatorname{dom}(f_j)\), and
\(\mathcal M\models\varphi(\bar a)\) iff
\(\mathcal N\models\varphi(f_j(\bar a))=\varphi(h(\bar a))\). With
\(\varphi\) equal to \(v_1=v_2\) this makes \(h\) injective; with
\(R(\bar v)\), \(F(\bar v)=w\) and \(c=w\) it shows that \(h\) preserves and
reflects relations and commutes with functions and constants. So \(h\) is an
isomorphism.
\end{proof}
\end{annexitem}

\begin{annexitem}{annex:p03-w03-p03-039-karp}{thm:p03-w03-p03-039-karp}{Karp's theorem}
\begin{proof}
(3)\(\Rightarrow\)(2) is immediate.

(2)\(\Rightarrow\)(1). Let \(P\) be a partial isomorphism system. We show,
by induction on formulas \(\varphi\) with finitely many free variables
(a class closed under taking constituents), that for every \(f\in P\) and
every assignment \(s\) into \(M\) whose values on the free variables of
\(\varphi\) lie in \(\operatorname{dom}(f)\),
\(\mathcal M,s\models\varphi\) iff \(\mathcal N,f\circ s\models\varphi\),
where \(f\circ s\) is any assignment into \(N\) agreeing with \(f(s(x))\) at
the free variables \(x\) of \(\varphi\). Atomic formulas: \(f\) is a partial
embedding. \(\neg\), \(\bigwedge\), \(\bigvee\): induction hypothesis with the
same \(f\). Let \(\varphi=\exists w\,\psi\). If
\(\mathcal M,s[w\mapsto c]\models\psi\), clause~1 gives \(g\in P\),
\(g\supseteq f\), with \(c\in\operatorname{dom}(g)\); the induction hypothesis
for \(g\) gives \(\mathcal N,(g\circ s)[w\mapsto g(c)]\models\psi\), and
\(g\circ s\) agrees with \(f\circ s\) on the free variables of \(\varphi\), so
\(\mathcal N,f\circ s\models\varphi\). If
\(\mathcal N,(f\circ s)[w\mapsto d]\models\psi\), clause~2 gives
\(g\supseteq f\) in \(P\) and \(c\) with \(g(c)=d\), and the induction
hypothesis for \(g\) gives \(\mathcal M,s[w\mapsto c]\models\psi\). The case
\(\forall w\) follows by the same two moves, or from the case \(\exists\)
applied to \(\neg\). Applying the claim to sentences and any \(f\in P\)
(\(P\ne\varnothing\)) gives \(\mathcal M\equiv_{\infty,\omega}\mathcal N\).

(4)\(\Rightarrow\)(3). Let \(\eta\) be a winning strategy for Player~II and
let \(P\) be the set of finite maps \(a_i\mapsto b_i\), \(i<n\), arising from
finite partial plays in which Player~II follows \(\eta\). The empty play
gives the empty map, so \(P\ne\varnothing\). A finite partial play following
\(\eta\) extends to a complete play following \(\eta\) (let Player~I continue
arbitrarily), which Player~II wins; the map of the partial play is the
restriction of a partial embedding to a finite set and so is a partial
embedding. Given such a map and \(a\in M\) (or \(b\in N\)), let Player~I play
\(a\) (or \(b\)) next and Player~II answer by \(\eta\); the new map is in
\(P\) and extends the old one, giving clauses~1 and~2. All members of \(P\)
have finite domain.

(1)\(\Rightarrow\)(4). Write \((\mathcal M,\bar a)\equiv_{\infty,\omega}(\mathcal N,\bar b)\)
when \(\mathcal M\models\varphi(\bar a)\Leftrightarrow\mathcal N\models\varphi(\bar b)\)
for every formula \(\varphi(\bar v)\) whose free variables are among
\(\bar v\). Claim: if this holds and \(c\in M\), there is \(d\in N\) with
\((\mathcal M,\bar a,c)\equiv_{\infty,\omega}(\mathcal N,\bar b,d)\); and
symmetrically for \(d\in N\). Suppose not. For each \(d\in N\) there is a
formula \(\varphi_d(\bar v,w)\) true at \((\bar a,c)\) in \(\mathcal M\) and
false at \((\bar b,d)\) in \(\mathcal N\), replacing a formula by its
negation if the difference goes the other way. Formulas form a proper class;
to choose, take for each \(d\) the set of such formulas of least
set-theoretic rank, and apply the axiom of choice to this set-indexed family
of nonempty sets. Then \(\mathcal M\models\exists w\bigwedge_{d\in N}\varphi_d(\bar a,w)\),
witnessed by \(c\), while \(\mathcal N\) satisfies the negation at \(\bar b\),
because every \(d\) falsifies \(\varphi_d\). The free variables of this
formula are among \(\bar v\), contradicting the hypothesis. The symmetric
claim is proved in the same way.

Player~II's strategy: after moves \(\bar a=(a_0,\ldots,a_{n-1})\),
\(\bar b=(b_0,\ldots,b_{n-1})\) with
\((\mathcal M,\bar a)\equiv_{\infty,\omega}(\mathcal N,\bar b)\), answer
Player~I's next move by an element given by the claim (chosen by a fixed
choice function on the nonempty sets of admissible answers). At the start
the condition is \(\mathcal M\equiv_{\infty,\omega}\mathcal N\). So the
condition holds at every finite stage. Any atomic formula concerns finitely
many \(a_i\), all played by some finite stage, at which the two tuples agree
on it; hence the map \(a_i\mapsto b_i\) is well defined (apply this to
\(v_i=v_j\)) and is a partial \(\tau\)-embedding. So the strategy wins.
\end{proof}
\end{annexitem}

\begin{annexitem}{annex:p03-w03-p03-039-approximant-basics}{thm:p03-w03-p03-039-approximant-basics}{Monotonicity and equivalence}
\begin{proof}
Write \(\bar a\sim_\alpha\bar b\) for
\((\mathcal M,\bar a)\sim_\alpha(\mathcal N,\bar b)\) when the structures are
clear.

\emph{Step 1: \(\sim_{\alpha+1}\) implies \(\sim_\alpha\), for all \(\alpha\).}
By induction on \(\alpha\), for all structures and tuples simultaneously.
For \(\alpha=0\): pick \(c\in M\) (structures are nonempty); the forth
condition gives \(d\) with \(\bar ac\sim_0\bar bd\), and the atomic formulas
not mentioning the last place show \(\bar a\sim_0\bar b\). For
\(\alpha=\gamma+1\): the forth and back conditions of \(\sim_{\gamma+2}\)
give \(\sim_{\gamma+1}\)-matches, which are \(\sim_\gamma\)-matches by the
induction hypothesis; that is \(\bar a\sim_{\gamma+1}\bar b\). For a limit
\(\alpha\): \(\sim_{\alpha+1}\) gives \(\sim_\alpha\)-matches, which are
\(\sim_\delta\)-matches for every \(\delta<\alpha\); so
\(\bar a\sim_{\delta+1}\bar b\) for every \(\delta<\alpha\), and by the
induction hypothesis \(\bar a\sim_\delta\bar b\) for every \(\delta<\alpha\),
which is \(\bar a\sim_\alpha\bar b\).

\emph{Step 2: monotonicity.} Fix \(\alpha\) and prove by induction on
\(\beta\ge\alpha\) that \(\sim_\beta\) implies \(\sim_\alpha\): trivial at
\(\beta=\alpha\); at \(\beta=\gamma+1>\alpha\), Step~1 and the induction
hypothesis for \(\gamma\ge\alpha\); at a limit \(\beta>\alpha\), the
definition.

\emph{Step 3: equivalence, permutation and deletion.} Each property is
proved by induction on \(\alpha\), for all tuples at once; the limit steps
follow from the definition, so only the zero and successor steps are given.
Reflexivity: \(\bar a\sim_0\bar a\); at a successor, match \(c\) with \(c\).
Symmetry: the definition is symmetric in the two sides. Transitivity across
\(\mathcal M,\mathcal N,\mathcal K\): agreement on atomic formulas is
transitive; at a successor, given \(c\in M\) take \(d\in N\) with
\(\bar ac\sim_\alpha\bar bd\), then \(e\in K\) with
\(\bar bd\sim_\alpha\bar ce\), and apply the induction hypothesis; the back
direction is symmetric. Permutation: for a permutation \(\pi\) of the places,
\(\bar a^\pi\) and \(\bar b^\pi\) satisfy the same atomic formulas as
\(\bar a\) and \(\bar b\) after renaming variables; at a successor,
\(\bar a^\pi c\) is \((\bar ac)^{\pi'}\) with \(\pi'\) fixing the last place.
Deletion of the same place \(i\) from both tuples: atomic formulas not
mentioning place \(i\) are among those already agreed on; at a successor,
\((\bar a\setminus a_i)c\) arises from \(\bar ac\) by deleting place \(i\),
and the induction hypothesis applies to the match of \(\bar ac\).
\end{proof}
\end{annexitem}

\begin{annexitem}{annex:p03-w03-p03-039-scott-formula-lemma}{thm:p03-w03-p03-039-scott-formula-lemma}{What the Scott formulas express}
\begin{proof}
By induction on \(\alpha\), for all tuples and all \(\mathcal N\). Take
\(\bar v=(v_1,\ldots,v_l)\) and \(w=v_{l+1}\). For \(\alpha=0\):
\(\mathcal N\models\Phi^{\mathcal M}_{\bar a,0}(\bar b)\) says that \(\bar b\)
satisfies every atomic formula true at \(\bar a\) and, through the negated
atomic conjuncts, no atomic formula false at \(\bar a\); that is
\(\bar a\sim_0\bar b\). For a limit \(\gamma\): both sides are conjunctions
over \(\alpha<\gamma\) of the corresponding statements. Successor
\(\alpha+1\): suppose \(\mathcal N\models\Phi^{\mathcal M}_{\bar a,\alpha+1}(\bar b)\).
For \(c\in M\), the conjunct \(\exists w\,\Phi^{\mathcal M}_{\bar a,c,\alpha}\)
yields \(d\in N\) with \(\mathcal N\models\Phi^{\mathcal M}_{\bar a,c,\alpha}(\bar b,d)\),
i.e. \(\bar ac\sim_\alpha\bar bd\) by the induction hypothesis. For
\(d\in N\), the conjunct \(\forall w\bigvee_{c\in M}\Phi^{\mathcal M}_{\bar a,c,\alpha}\)
yields \(c\in M\) with \(\mathcal N\models\Phi^{\mathcal M}_{\bar a,c,\alpha}(\bar b,d)\),
i.e. \(\bar ac\sim_\alpha\bar bd\). So \(\bar a\sim_{\alpha+1}\bar b\).
Conversely, if \(\bar a\sim_{\alpha+1}\bar b\), the forth condition and the
induction hypothesis make every conjunct
\(\exists w\,\Phi^{\mathcal M}_{\bar a,c,\alpha}(\bar b,w)\) true, and the
back condition makes \(\bigvee_{c}\Phi^{\mathcal M}_{\bar a,c,\alpha}(\bar b,d)\)
true for every \(d\).

Quantifier rank, by induction on \(\alpha\): the conjuncts at stage \(0\) are
quantifier free; at \(\alpha+1\) each displayed part is a quantifier applied
to a formula of rank at most \(\alpha\), so has rank at most \(\alpha+1\);
at a limit the supremum of ranks at most \(\beta<\alpha\) is at most
\(\alpha\).
\end{proof}
\end{annexitem}

\begin{annexitem}{annex:p03-w03-p03-039-rank-characterization}{thm:p03-w03-p03-039-rank-characterization}{Back-and-forth rank equals quantifier rank}
\begin{proof}
(\(\Rightarrow\)) We prove by induction on formulas \(\varphi(\bar v)\) with
free variables among a finite list \(\bar v\): for every \(\alpha\ge
\operatorname{qr}(\varphi)\) and all \(\bar a\sim_\alpha\bar b\) of length
\(\lvert\bar v\rvert\), \(\mathcal M\models\varphi(\bar a)\) iff
\(\mathcal N\models\varphi(\bar b)\). Atomic: \(\sim_\alpha\) implies
\(\sim_0\) by \cref{thm:p03-w03-p03-039-approximant-basics}. \(\neg\),
\(\bigwedge\), \(\bigvee\): every constituent has rank at most \(\alpha\).
Let \(\varphi=\exists w\,\psi\) with \(\gamma=\operatorname{qr}(\psi)\),
so \(\gamma+1\le\alpha\) and \(\bar a\sim_{\gamma+1}\bar b\) by monotonicity.
If \(w\) is one of the variables of \(\bar v\), delete that place from
\(\bar v,\bar a,\bar b\); deletion preserves \(\sim_{\gamma+1}\)
(\cref{thm:p03-w03-p03-039-approximant-basics}) and \(\varphi\) does not
have \(w\) free. Now if \(\mathcal M\models\psi(\bar a,c)\), the forth
condition gives \(d\) with \(\bar ac\sim_\gamma\bar bd\), and the induction
hypothesis for \(\psi(\bar v,w)\) at \(\gamma\) gives
\(\mathcal N\models\psi(\bar b,d)\); the converse uses the back condition.
The case \(\forall w\,\psi\) is the same with the two conditions used to
transfer counterexamples.

(\(\Leftarrow\)) Suppose \(\bar a\) and \(\bar b\) agree on all formulas of
rank at most \(\alpha\) with free variables among \(\bar v\). The formula
\(\Phi^{\mathcal M}_{\bar a,\alpha}(\bar v)\) has free variables among
\(\bar v\) and rank at most \(\alpha\)
(\cref{thm:p03-w03-p03-039-scott-formula-lemma}); it holds at \(\bar a\) in
\(\mathcal M\) because \(\bar a\sim_\alpha\bar a\)
(\cref{thm:p03-w03-p03-039-approximant-basics}) and by
\cref{thm:p03-w03-p03-039-scott-formula-lemma}. So it holds at \(\bar b\) in
\(\mathcal N\), and the same lemma gives \(\bar a\sim_\alpha\bar b\). The
unpointed statement is the case of empty tuples.
\end{proof}
\end{annexitem}

\begin{annexitem}{annex:p03-w03-p03-039-scott-rank}{thm:p03-w03-p03-039-scott-rank}{The Scott rank of an infinite structure}
\begin{proof}
By \cref{thm:p03-w03-p03-039-approximant-basics}, \(\Gamma_\beta\subseteq\Gamma_\alpha\)
for \(\alpha\le\beta\).

Claim~1: if \(\Gamma_\alpha=\Gamma_{\alpha+1}\), then \(\Gamma_\beta=\Gamma_\alpha\)
for all \(\beta\ge\alpha\). Induction on \(\beta\). At a limit \(\beta\),
\(\Gamma_\beta=\bigcap_{\gamma<\beta}\Gamma_\gamma\); the terms with
\(\gamma\ge\alpha\) equal \(\Gamma_\alpha\) and those with \(\gamma<\alpha\)
contain it, so the intersection is \(\Gamma_\alpha\). At \(\beta+1\) with
\(\beta\ge\alpha\): let \((\bar a,\bar b)\in\Gamma_\beta=\Gamma_\alpha=\Gamma_{\alpha+1}\).
For \(c\in M\) there is \(d\) with \((\bar ac,\bar bd)\in\Gamma_\alpha=\Gamma_\beta\),
and symmetrically, so \((\bar a,\bar b)\in\Gamma_{\beta+1}\); with
\(\Gamma_{\beta+1}\subseteq\Gamma_\beta\) this gives equality.

Claim~2: there is \(\alpha<\lvert M\rvert^+\) with \(\Gamma_\alpha=\Gamma_{\alpha+1}\).
Otherwise choose, for each \(\alpha<\lvert M\rvert^+\), a pair
\((\bar a_\alpha,\bar b_\alpha)\in\Gamma_\alpha\setminus\Gamma_{\alpha+1}\).
If \(\alpha<\beta\), then \((\bar a_\beta,\bar b_\beta)\in\Gamma_\beta\subseteq\Gamma_{\alpha+1}\),
so the pairs differ. This injects \(\lvert M\rvert^+\) into the set of pairs
of finite tuples from \(M\), which has cardinality \(\lvert M\rvert\) because
\(M\) is infinite; impossible.

For \(\alpha\) as in Claim~2 and any \(\beta\), Claim~1 gives
\(\Gamma_\alpha=\Gamma_\beta\) when \(\beta\ge\alpha\), and
\(\Gamma_\alpha\subseteq\Gamma_\beta\) when \(\beta<\alpha\); so
\(\sim_\alpha\) implies \(\sim_\beta\) on \(M\). Conversely, if \(\sim_\alpha\)
implies every \(\sim_\beta\), then \(\Gamma_\alpha\subseteq\Gamma_{\alpha+1}\subseteq\Gamma_\alpha\).
Hence the two descriptions of the least such \(\alpha\) agree, and it is
below \(\lvert M\rvert^+\).
\end{proof}
\end{annexitem}

\begin{annexitem}{annex:p03-w03-p03-039-tuple-rank-bound}{thm:p03-w03-p03-039-tuple-rank-bound}{Tuple ranks are bounded by the Scott rank}
\begin{proof}
Let \(\alpha=\operatorname{sr}(\mathcal M)\). By
\cref{thm:p03-w03-p03-039-scott-rank}, \((\mathcal M,\bar a)\sim_\alpha(\mathcal M,\bar b)\)
implies \((\mathcal M,\bar a)\sim_{\alpha+1}(\mathcal M,\bar b)\) for all
\(\bar b\). So \(\alpha\) satisfies the condition defining \(r(\bar a)\),
and \(r(\bar a)\le\alpha\).
\end{proof}
\end{annexitem}

\begin{annexitem}{annex:p03-w03-p03-039-sr-counterexample}{thm:p03-w03-p03-039-sr-counterexample}{A structure whose Scott rank exceeds every tuple rank}
\begin{proof}
Write \(A=\{a_1,a_2,a_3\}\), \(B=\{b_1,b_2\}\), and \(\bar c\sim_\beta\bar d\)
for \((\mathcal M,\bar c)\sim_\beta(\mathcal M,\bar d)\). For a tuple
\(\bar c\), a \emph{block} of \(\bar c\) is the set of entries of \(\bar c\)
lying in one of the classes \(A\), \(B\). If \(\bar c\sim_0\bar d\), the two
tuples have the same pattern of equalities, of \(E\) and of \(P\), so
\(c_i\mapsto d_i\) is a well-defined bijection between the sets of entries
that carries each block of \(\bar c\) onto a block of \(\bar d\) of the same
size. Call \(\bar c\sim_0\bar d\) \emph{class-matched} if \(c_i\in A\)
exactly when \(d_i\in A\), for every \(i\).

Step~1: if \(\bar c\sim_0\bar d\) is class-matched, then \(\bar c\sim_\beta\bar d\)
for every \(\beta\). The map \(c_i\mapsto d_i\) is injective and sends entries
in \(A\), in \(B\) and in \(\omega\) into the same set; extend its restriction
to each of these three sets to a permutation of that set (a finite injective
partial map on a set extends to a permutation of it). The union \(\sigma\) of
the three permutations preserves \(P\) and \(E\), so it is an automorphism
with \(\sigma(\bar c)=\bar d\). An automorphism \(\sigma\) gives
\(\bar e\sim_\beta\sigma(\bar e)\) for every tuple \(\bar e\) and every
\(\beta\), by induction on \(\beta\): at \(0\) because \(\sigma\) preserves
atomic formulas, at a successor by answering \(e\) with \(\sigma(e)\) and
\(f\) with \(\sigma^{-1}(f)\), and at limits trivially.

Step~2: if \(\bar c\sim_1\bar d\), every block of \(\bar c\) with at least
two elements is carried into the same class. Suppose a block \(X\) of
\(\bar c\) with \(\lvert X\rvert=m\ge2\) is carried onto a block \(Y\) of
\(\bar d\) in the other class. As \(\lvert B\rvert=2\), \(m=2\), and one of
\(X,Y\) is all of \(B\) while the other is a two-element subset of \(A\).
If \(X\subseteq A\), let \(e\) be the element of \(A\setminus X\); a
\(\sim_0\)-answer \(f\) to \(e\) would satisfy \(f\mathrel E y\) and
\(f\notin Y\) for \(y\in Y=B\), and no such \(f\) exists. If \(Y\subseteq A\),
the element of \(A\setminus Y\) has no answer in the back direction for the
same reason. Either way \(\bar c\not\sim_1\bar d\).

Step~3: if \(\bar c\sim_2\bar d\), then \(\bar c\sim_0\bar d\) is
class-matched. By Step~2 (and \(\sim_2\Rightarrow\sim_1\),
\cref{thm:p03-w03-p03-039-approximant-basics}) only a one-element block
\(\{c_i\}\) can be carried into the other class. Let \(e\ne c_i\) be in the
class of \(c_i\); it is not an entry of \(\bar c\). The forth clause of
\(\sim_2\) gives \(f\) with \(\bar ce\sim_1\bar df\). In \(\bar ce\) the
block of \(c_i\) is \(\{c_i,e\}\), and it is carried to the block of \(d_i\)
in \(\bar df\), which lies in the other class; this contradicts Step~2.

Step~4: \(\operatorname{sr}(\mathcal M)=2\). By Steps~3 and~1, \(\sim_2\)
implies every \(\sim_\beta\); so \(\Gamma_2=\Gamma_3\) and
\(\operatorname{sr}(\mathcal M)\le2\)
(\cref{thm:p03-w03-p03-039-scott-rank}). Next, \(a_1\sim_1b_1\): an answer
to \(a_1\) is \(b_1\); to \(a_2\) or \(a_3\) it is \(b_2\); to \(b_1\) or
\(b_2\) it is \(a_2\); to \(n\in\omega\) it is \(n\); and the back direction
answers \(b_1\), \(b_2\), \(a_i\), \(n\) by \(a_1\), \(a_2\), \(b_1\), \(n\).
In each case the two pairs have the same pattern of equality, \(E\) and
\(P\). By Step~3, \(a_1\not\sim_2b_1\), since \(a_1\in A\) and
\(b_1\notin A\). So \(\Gamma_1\ne\Gamma_2\). If \(\Gamma_0=\Gamma_1\) or
\(\Gamma_1=\Gamma_2\), Claim~1 in the proof of
\cref{thm:p03-w03-p03-039-scott-rank} would give \(\Gamma_1=\Gamma_2\); so
the least \(\alpha\) with \(\Gamma_\alpha=\Gamma_{\alpha+1}\) is \(2\).

Step~5: \(r(\bar c)\le1\) for every \(\bar c\). Case~(i): \(\bar c\) has a
block with at least two elements. Let \(\bar c\sim_1\bar d\). By Step~2
that block is carried into its own class. If \(\bar c\) has a second block,
it lies in the other class, and its image lies in a block of \(\bar d\)
different from the image of the first, hence in the other class as well.
So \(\bar c\sim_0\bar d\) is class-matched and \(\bar c\sim_2\bar d\) by
Step~1: the condition defining \(r(\bar c)\) holds at \(1\). Case~(ii): every
block of \(\bar c\) has at most one element. Let \(\bar c\sim_0\bar d\); then
every block of \(\bar d\) has at most one element too. Given \(e\in M\): if
\(e=c_i\), answer \(d_i\); if \(e\in\omega\) is new, answer a new element of
\(\omega\); if \(e\notin\omega\) is new and lies in the class of an entry
\(c_i\), answer an element of the class of \(d_i\) other than \(d_i\), which
is not an entry of \(\bar d\) because that block of \(\bar d\) is
\(\{d_i\}\); if \(e\notin\omega\) is new and its class contains no entry of
\(\bar c\), then \(\bar c\) has at most one block, so does \(\bar d\), and
an element of a class containing no entry of \(\bar d\) is an answer. In
each case the extended tuples have the same pattern of equality, \(E\) and
\(P\). The back direction is the same argument with the tuples exchanged.
So \(\bar c\sim_1\bar d\), and the condition defining \(r(\bar c)\) holds at
\(0\).

Hence \(r(a_1)=0\) by Case~(ii), and \(r(a_1,a_2)=1\): it is at most \(1\)
by Case~(i), and it is not \(0\) because \((a_1,a_2)\sim_0(b_1,b_2)\) while
\((a_1,a_2)\not\sim_1(b_1,b_2)\) by Step~2. So
\(\sup_{\bar c}r(\bar c)=1\), \(\operatorname{SR}(\mathcal M)=2\) and
\(\operatorname{sr}(\mathcal M)=2\).
\end{proof}
\end{annexitem}

\begin{annexitem}{annex:p03-w03-p03-039-scott}{thm:p03-w03-p03-039-scott}{Scott's isomorphism theorem}
\begin{proof}
Let \(\alpha=\operatorname{sr}(\mathcal M)\).

\(\mathcal M\models\Phi^{\mathcal M}\): \(\Phi^{\mathcal M}_{\varnothing,\alpha}\)
holds in \(\mathcal M\) by reflexivity of \(\sim_\alpha\) and
\cref{thm:p03-w03-p03-039-scott-formula-lemma}. If
\(\mathcal M\models\Phi^{\mathcal M}_{\bar a,\alpha}(\bar b)\), then
\((\mathcal M,\bar a)\sim_\alpha(\mathcal M,\bar b)\), hence
\(\sim_{\alpha+1}\) by \cref{thm:p03-w03-p03-039-scott-rank}, hence
\(\mathcal M\models\Phi^{\mathcal M}_{\bar a,\alpha+1}(\bar b)\). So if
\(\mathcal N\equiv_{\infty,\omega}\mathcal M\), then
\(\mathcal N\models\Phi^{\mathcal M}\).

Conversely let \(\mathcal N\models\Phi^{\mathcal M}\). Let \(P\) be the set of
maps \(f\) from a finite subset of \(M\) into \(N\) such that, for some
enumeration \(\bar a\) of \(\operatorname{dom}(f)\),
\(\mathcal N\models\Phi^{\mathcal M}_{\bar a,\alpha}(f(\bar a))\). By
\cref{thm:p03-w03-p03-039-scott-formula-lemma} this says
\((\mathcal M,\bar a)\sim_\alpha(\mathcal N,f(\bar a))\), which by
\cref{thm:p03-w03-p03-039-approximant-basics} does not depend on the
enumeration and implies \(\sim_0\); so every \(f\in P\) is a partial
\(\tau\)-embedding. The empty map is in \(P\) because \(\mathcal N\models\Phi^{\mathcal M}_{\varnothing,\alpha}\).

Let \(f\in P\) with \(\bar b=f(\bar a)\). The conjunct of \(\Phi^{\mathcal M}\)
for \(\bar a\), instantiated at \(\bar b\), gives
\(\mathcal N\models\Phi^{\mathcal M}_{\bar a,\alpha+1}(\bar b)\), i.e.
\((\mathcal M,\bar a)\sim_{\alpha+1}(\mathcal N,\bar b)\). Given \(c\in M\),
there is \(d\in N\) with \((\mathcal M,\bar ac)\sim_\alpha(\mathcal N,\bar bd)\).
If \(c=a_i\), agreement on \(v_i=w\) forces \(d=b_i\); so
\(g=f\cup\{(c,d)\}\) is a function, and \(g\in P\) with enumeration
\(\bar ac\). Given \(d\in N\), the back condition gives \(c\in M\) with
\((\mathcal M,\bar ac)\sim_\alpha(\mathcal N,\bar bd)\), and the same
reasoning puts \(f\cup\{(c,d)\}\) in \(P\). So \(P\) is a partial
isomorphism system, \(\mathcal M\cong_p\mathcal N\), and
\(\mathcal M\equiv_{\infty,\omega}\mathcal N\) by
\cref{thm:p03-w03-p03-039-karp}.

If \(\tau\) and \(\mathcal M\) are countable, \(\Phi^{\mathcal M}\in\mathcal L_{\omega_1,\omega}\)
by \cref{def:p03-w03-p03-039-scott-sentence}, and if moreover \(\mathcal N\)
is countable and \(\mathcal N\models\Phi^{\mathcal M}\), the system \(P\)
gives \(\mathcal M\cong\mathcal N\) by
\cref{thm:p03-w03-p03-039-countable-back-forth}.
\end{proof}
\end{annexitem}

\begin{annexitem}{annex:p03-w03-p03-039-scott-complete}{thm:p03-w03-p03-039-scott-complete}{Scott sentences are complete}
\begin{proof}
\(\Phi^{\mathcal M}\) is satisfiable, since \(\mathcal M\models\Phi^{\mathcal M}\)
(proof of \cref{thm:p03-w03-p03-039-scott}). Let \(\psi\) be any
\(\mathcal L_{\infty,\omega}\)-sentence. If \(\mathcal M\models\psi\), every
model \(\mathcal N\) of \(\Phi^{\mathcal M}\) has
\(\mathcal N\equiv_{\infty,\omega}\mathcal M\) by
\cref{thm:p03-w03-p03-039-scott} and so satisfies \(\psi\); hence
\(\Phi^{\mathcal M}\models\psi\). Otherwise \(\mathcal M\models\neg\psi\)
and in the same way \(\Phi^{\mathcal M}\models\neg\psi\).
\end{proof}
\end{annexitem}

\begin{annexitem}{annex:p03-w03-p03-039-consistency-lemma}{thm:p03-w03-p03-039-consistency-lemma}{Equality and modus ponens inside a consistency property}
\begin{proof}
(1) By C7(c) with \(t=c\) there are \(d\in C\) and \(\mu_1\in\Sigma\) with
\(\sigma\cup\{d=c\}\subseteq\mu_1\). By C7(a) applied to \(\mu_1\) there is
\(\mu_2\in\Sigma\) containing \(\mu_1\cup\{c=d\}\). Apply C7(b) to
\(\mu_2\) with the constant \(c\), the closed term \(d\) and the formula
\(\varphi(v)\) equal to \(v=c\): its premises \(c=d\) and
\(\varphi(d)=(d=c)\) are in \(\mu_2\). C7(b) gives \(\mu\in\Sigma\) containing
\(\mu_2\cup\{\varphi(c)\}=\mu_2\cup\{c=c\}\), and \(\sigma\subseteq\mu\).

(2) By C7(a) applied to \(d=e\in\sigma\), some \(\mu_1\in\Sigma\) contains
\(\sigma\cup\{e=d\}\). Apply C7(b) to \(\mu_1\) with the constant \(e\),
the term \(d\) and \(\varphi(v)\) equal to \(c=v\): \(e=d\in\mu_1\) and
\(\varphi(d)=(c=d)\in\mu_1\), so some \(\mu\in\Sigma\) contains
\(\mu_1\cup\{c=e\}\supseteq\sigma\cup\{c=e\}\). By C0,
\(\sigma\cup\{c=e\}\in\Sigma\).

(3) \(\varphi\to\psi\) is \(\bigvee\{\neg\varphi,\psi\}\). By C4 there is
\(\mu\in\Sigma\) containing \(\sigma\) and one of \(\neg\varphi\), \(\psi\).
If it contained \(\neg\varphi\), it would contain both \(\varphi\) and
\(\neg\varphi\), against C1. So \(\psi\in\mu\).
\end{proof}
\end{annexitem}

\begin{annexitem}{annex:p03-w03-p03-039-model-existence}{thm:p03-w03-p03-039-model-existence}{Model existence theorem}
\begin{proof}
\emph{The sentences used.} Let \(\Lambda\) be the least set of
\(\mathcal L_{\omega_1,\omega}(\tau)\)-sentences such that: \(\sigma\subseteq\Lambda\);
every atomic sentence and every negated atomic sentence of \(\tau\) is in
\(\Lambda\); every subformula of a member of \(\Lambda\) that is a sentence is
in \(\Lambda\); if \(\varphi(v)\) is a subformula of a member of \(\Lambda\)
with at most \(v\) free and \(c\in C\), then \(\varphi(c)\in\Lambda\); and if
\(\neg\varphi\in\Lambda\) then \({\sim}\varphi\in\Lambda\). The atomic sentences
are countable because \(\tau\) is, each sentence has countably many
subformulas, and \(C\) is countable; building \(\Lambda\) in \(\omega\)
countable stages shows that \(\Lambda\) is countable. List \(\Lambda\) as
\(\varphi_0,\varphi_1,\ldots\) with every member occurring infinitely often,
and list the closed terms as \(t_0,t_1,\ldots\).

\emph{The chain.} Put \(\sigma_0=\sigma\). Given \(\sigma_n\in\Sigma\) with
\(\sigma_n\subseteq\Lambda\): if \(\sigma_n\cup\{\varphi_n\}\notin\Sigma\)
put \(\sigma_n'=\sigma_n\). Otherwise, if \(\varphi_n\) is
\(\bigvee_{\psi\in X}\psi\), C4 applied to \(\sigma_n\cup\{\varphi_n\}\)
gives \(\psi\in X\) and, by C0, \(\sigma_n'=\sigma_n\cup\{\varphi_n,\psi\}\in\Sigma\);
if \(\varphi_n\) is \(\exists v\,\psi(v)\), C6 gives \(c\in C\) and
\(\sigma_n'=\sigma_n\cup\{\varphi_n,\psi(c)\}\in\Sigma\); otherwise
\(\sigma_n'=\sigma_n\cup\{\varphi_n\}\). Then C7(c) gives \(e\in C\) with
\(\sigma_{n+1}=\sigma_n'\cup\{e=t_n\}\in\Sigma\) (by C0). All added sentences
are in \(\Lambda\). Let \(\Gamma=\bigcup_n\sigma_n\subseteq\Lambda\).

\emph{Closure facts.} (a) A finite subset of \(\Gamma\) lies in some
\(\sigma_n\); with C1, no \(\varphi\) and \(\neg\varphi\) are both in
\(\Gamma\). (b) If \(\theta\in\Lambda\) and \(\sigma_n\cup\{\theta\}\in\Sigma\)
for all large \(n\), then \(\theta\in\Gamma\): take a large \(n\) with
\(\varphi_n=\theta\). (c) Whenever C2, C3, C5 or C7(a),(b), or
\cref{thm:p03-w03-p03-039-consistency-lemma}, applied to premises in
\(\sigma_n\), yields \(\mu\supseteq\sigma_n\cup\{\theta\}\), C0 gives
\(\sigma_n\cup\{\theta\}\in\Sigma\); premises in \(\Gamma\) lie in all large
\(\sigma_n\), so by (b) the conclusion \(\theta\) is in \(\Gamma\) whenever
\(\theta\in\Lambda\). Thus: \(\neg\varphi\in\Gamma\) gives
\({\sim}\varphi\in\Gamma\); \(\bigwedge X\in\Gamma\) gives \(X\subseteq\Gamma\);
\(\forall v\,\varphi\in\Gamma\) gives \(\varphi(c)\in\Gamma\) for all
\(c\in C\); \(c=c\in\Gamma\); \(c=d\in\Gamma\) gives \(d=c\in\Gamma\);
\(c=d,d=e\in\Gamma\) give \(c=e\in\Gamma\); and \(c=t,\varphi(t)\in\Gamma\)
give \(\varphi(c)\in\Gamma\) whenever \(\varphi(c)\) is atomic or negated
atomic. (d) If \(\bigvee X\in\Gamma\), take a large \(n\) with
\(\varphi_n=\bigvee X\in\sigma_n\); then \(\sigma_n\cup\{\varphi_n\}=\sigma_n\in\Sigma\)
and the construction put some \(\psi\in X\) into \(\Gamma\). Likewise
\(\exists v\,\psi\in\Gamma\) gives \(\psi(c)\in\Gamma\) for some \(c\in C\).
(e) For every closed term \(t\) there is \(e\in C\) with \(e=t\in\Gamma\).

\emph{The structure.} For \(c,d\in C\) let \(c\approx d\) iff \(c=d\in\Gamma\).
By (c), \(\approx\) is an equivalence relation. Let \(M=\{[c]:c\in C\}\),
countable and nonempty. For a function symbol \(f\) (constants being the
\(0\)-ary case) and \(c_1,\ldots,c_n\in C\), (e) gives \(d\) with
\(d=f(c_1,\ldots,c_n)\in\Gamma\); put \(f^{\mathcal M}([c_1],\ldots,[c_n])=[d]\).
This is well defined: if \(d'=f(c_1',\ldots,c_n')\in\Gamma\) with
\(c_i\approx c_i'\), replace each \(c_i'\) by \(c_i\) using \(c_i=c_i'\in\Gamma\)
and the last item of (c), obtaining \(d'=f(c_1,\ldots,c_n)\in\Gamma\); then
the same item with \(d=f(\bar c)\) and \(\varphi(v)\) equal to \(d'=v\) gives
\(d'=d\in\Gamma\). For a relation symbol \(R\), let
\(R^{\mathcal M}([c_1],\ldots,[c_n])\) iff \(R(c_1,\ldots,c_n)\in\Gamma\),
well defined by the same substitution item.

Claim: if \(t\) is a closed term and \(c=t\in\Gamma\), then
\(t^{\mathcal M}=[c]\). By induction on \(t=f(s_1,\ldots,s_m)\) (constants
included with \(m=0\)): by (e) take \(d_i\) with \(d_i=s_i\in\Gamma\), so
\(s_i^{\mathcal M}=[d_i]\) by the induction hypothesis; replacing \(s_i\) by
\(d_i\) one place at a time in \(c=f(s_1,\ldots,s_m)\) by the substitution
item of (c) (all intermediate sentences are atomic) gives
\(c=f(d_1,\ldots,d_m)\in\Gamma\), so \(t^{\mathcal M}=f^{\mathcal M}([d_1],\ldots,[d_m])=[c]\).

\emph{Truth.} Define a depth \(\delta\) on formulas: \(0\) on atomic formulas,
\(\delta(\neg\varphi)=\delta(\varphi)+1\),
\(\delta(\bigwedge X)=\delta(\bigvee X)=\sup\{\delta(\psi)+1:\psi\in X\}\),
\(\delta(\exists v\,\varphi)=\delta(\forall v\,\varphi)=\delta(\varphi)+1\).
Substituting constants does not change \(\delta\), and by induction
\(\delta({\sim}\varphi)\le\delta(\varphi)+1\). We prove by induction on
\(\delta(\varphi)\) that \(\varphi\in\Gamma\) implies \(\mathcal M\models\varphi\).

Atomic \(t_1=t_2\): take \(c_i=t_i\in\Gamma\) by (e). Substituting \(c_1\) for
\(t_1\) and then \(c_2\) for \(t_2\) (substitution item of (c)) gives
\(c_1=c_2\in\Gamma\), so \(t_1^{\mathcal M}=[c_1]=[c_2]=t_2^{\mathcal M}\) by the
claim. Atomic \(R(t_1,\ldots,t_n)\): in the same way \(R(c_1,\ldots,c_n)\in\Gamma\)
and \(\mathcal M\models R(\bar t)\). \(\bigwedge X\): (c) and the induction
hypothesis. \(\bigvee X\) and \(\exists v\,\psi\): (d) and the induction
hypothesis. \(\forall v\,\psi\): (c) gives \(\psi(c)\in\Gamma\) for every
\(c\in C\), every element of \(M\) is some \([c]\), and \([c]\) is the
value of \(c\) by the claim.

Negation \(\neg\psi\in\Gamma\); then \({\sim}\psi\in\Gamma\) by (c). If
\(\psi\) is \(t_1=t_2\): substitution as above turns \(\neg\,t_1=t_2\) into
\(\neg\,c_1=c_2\in\Gamma\); were \(\mathcal M\models t_1=t_2\), the claim
would give \([c_1]=[c_2]\), i.e. \(c_1=c_2\in\Gamma\), against (a). If
\(\psi\) is \(R(\bar t)\): the same substitutions give \(\neg R(\bar c)\in\Gamma\), so
\(R(\bar c)\notin\Gamma\) by (a), and \(\mathcal M\not\models R(\bar t)\).
If \(\psi=\neg\theta\): \({\sim}\psi=\theta\in\Gamma\) with
\(\delta(\theta)<\delta(\neg\neg\theta)\). If \(\psi=\bigwedge X\):
\(\bigvee_{\theta\in X}{\sim}\theta\in\Gamma\), so by (d) some
\({\sim}\theta\in\Gamma\); since
\(\delta({\sim}\theta)\le\delta(\theta)+1\le\delta(\psi)<\delta(\neg\psi)\),
the induction hypothesis and \cref{thm:p03-w03-p03-039-formal-negation}
give \(\mathcal M\models\neg\theta\), so \(\mathcal M\models\neg\psi\). If
\(\psi=\bigvee X\): every \({\sim}\theta\in\Gamma\), and in the same way
\(\mathcal M\models\neg\theta\) for all \(\theta\in X\). If
\(\psi=\exists v\,\theta\): \(\forall v\,{\sim}\theta\in\Gamma\), so
\({\sim}\theta(c)\in\Gamma\) for every \(c\) (it is in \(\Lambda\) by the
substitution clause), and the induction hypothesis gives
\(\mathcal M\models\neg\theta(c)\) for all \(c\), so
\(\mathcal M\models\neg\exists v\,\theta\). If \(\psi=\forall v\,\theta\):
\(\exists v\,{\sim}\theta\in\Gamma\) gives some \({\sim}\theta(c)\in\Gamma\)
and \(\mathcal M\models\neg\theta(c)\).

Since \(\sigma\subseteq\Gamma\), \(\mathcal M\models\sigma\), and \(\mathcal M\)
is countable.
\end{proof}
\end{annexitem}

\begin{annexitem}{annex:p03-w03-p03-039-extended-model-existence}{thm:p03-w03-p03-039-extended-model-existence}{Extended model existence theorem}
\begin{proof}
Since \(\Sigma\ne\varnothing\), C0 gives \(\varnothing\in\Sigma\). If
\(T=\varnothing\), \cref{thm:p03-w03-p03-039-model-existence} with
\(\sigma=\varnothing\) gives a countable structure, which models \(T\).
Otherwise list \(T\) as \(\theta_0,\theta_1,\ldots\) (with repetitions if
\(T\) is finite), start from \(\sigma_0=\varnothing\), add \(T\) to the
generators of \(\Lambda\), and in the construction of
\cref{thm:p03-w03-p03-039-model-existence} replace \(\sigma_{n+1}\) by
\(\sigma_{n+1}\cup\{\theta_n\}\), which lies in \(\Sigma\) by hypothesis.
All facts used in that proof only concern membership of each stage in
\(\Sigma\) and of each added sentence in \(\Lambda\), so they continue to
hold. Now \(T\subseteq\Gamma\), and the canonical structure is a countable
model of \(T\).
\end{proof}
\end{annexitem}

\begin{annexitem}{annex:p03-w03-p03-039-subformula-count}{thm:p03-w03-p03-039-subformula-count}{Subformula counts and regularity}
\begin{proof}
(1) Let \(\operatorname{Sub}(\varphi)\) be the set of subformulas. It is
computed by recursion: \(\operatorname{Sub}(\varphi)=\{\varphi\}\) for atomic
\(\varphi\);
\(\operatorname{Sub}(\neg\varphi)=\{\neg\varphi\}\cup\operatorname{Sub}(\varphi)\),
and likewise for \(\exists v\,\varphi\) and \(\forall v\,\varphi\); and
\(\operatorname{Sub}(\bigwedge X)=\{\bigwedge X\}\cup\bigcup_{\psi\in X}\operatorname{Sub}(\psi)\),
likewise for \(\bigvee X\). By induction on formulas,
\(\lvert\operatorname{Sub}(\varphi)\rvert<\kappa\): the atomic case gives
\(1\); the unary cases add one element to a set of size less than \(\kappa\);
and for \(\bigwedge X\) with \(\lvert X\rvert<\kappa\) the union is over fewer
than \(\kappa\) sets each of size less than \(\kappa\), which has size less
than \(\kappa\) because \(\kappa\) is regular (for \(\kappa=\omega\) this is
the fact that a finite union of finite sets is finite).

(2) Let \(\lambda=\operatorname{cf}(\kappa)<\kappa\) and let
\((\kappa_i:i<\lambda)\) be cardinals below \(\kappa\) with supremum
\(\kappa\). For \(\xi<\kappa\) let \(\chi_\xi\) be the sentence
\(\exists v_\xi\,v_\xi=v_\xi\); these sentences are pairwise distinct
because they use distinct variables, and each uses a variable \(v_\xi\) with
\(\xi<\kappa\). Let \(\psi_i=\bigwedge_{\xi<\kappa_i}\chi_\xi\), a conjunction
of \(\kappa_i<\kappa\) sentences, and \(\varphi=\bigwedge_{i<\lambda}\psi_i\),
a conjunction of \(\lambda<\kappa\) sentences. Then
\(\varphi\in\mathcal L_{\kappa,\omega}\), and every \(\chi_\xi\) with
\(\xi<\kappa\) is a subformula of \(\varphi\), because \(\xi<\kappa_i\) for
some \(i\). So \(\varphi\) has at least \(\kappa\) subformulas, and it has at
most \(\kappa\): the \(\chi_\xi\), their subformulas \(v_\xi=v_\xi\), the
\(\psi_i\) and \(\varphi\).
\end{proof}
\end{annexitem}

\begin{annexitem}{annex:p03-w03-p03-039-ultraproducts}{thm:p03-w03-p03-039-ultraproducts}{Ultraproducts and \(\mathcal L_{\omega_1,\omega}\)}
\begin{proof}
The interpretations in \cref{def:p03-w03-p03-039-ultraproduct} do not depend
on the representatives: if \(f_k\sim_Ug_k\) for \(k\le n\), the set where all
of them agree is a finite intersection of members of \(U\), hence in \(U\),
and the sets defining \(F\) and \(R\) for the \(g_k\) contain its
intersection with the sets for the \(f_k\). By induction on terms,
\(t([f_1],\dots,[f_n])=[i\mapsto t^{\mathcal M_i}(f_1(i),\dots,f_n(i))]\).

(1) Each \(\mathcal M_i\) satisfies \(\varphi\) because each \(m\in\omega\) is
the value of \(c_m\). Let \(g\) be the identity map of \(\omega\), viewed as
an element of \(\prod_i M_i\). For each \(n\),
\([g]=c_n^{\prod}=[i\mapsto n]\) would need \(\{i:i=n\}=\{n\}\in U\), which
fails because \(U\) is nonprincipal. So the element \([g]\) satisfies no
disjunct \(v=c_n\), and the ultraproduct satisfies \(\neg\varphi\).

(2) Write \(A_\psi(\bar f)=\{i:\mathcal M_i\models\psi(\bar f(i))\}\). The
proof is by induction on \(\varphi\); all subformulas of a formula with
finitely many free variables again have finitely many
(\cref{thm:p03-w03-p03-039-nadel-free-variables}), so the induction stays
among such formulas. Atomic formulas: by the term identity above and the
definition of the relations, including equality, which is the relation
\(\sim_U\). Negation: \(A_{\neg\psi}=I\setminus A_\psi\), and exactly one of a
set and its complement lies in the ultrafilter \(U\). Countable conjunction
\(\bigwedge_{n}\psi_n\): \(A=\bigcap_nA_{\psi_n}\). If every \(A_{\psi_n}\) is
in \(U\), so is \(A\) by \(\sigma\)-completeness; if \(A\in U\), each
\(A_{\psi_n}\supseteq A\) is in \(U\). With the induction hypothesis this is
the claim. Countable disjunction \(\bigvee_n\psi_n\): \(A=\bigcup_nA_{\psi_n}\).
If some \(A_{\psi_n}\in U\) then \(A\in U\). If \(A\in U\) and no
\(A_{\psi_n}\) is in \(U\), every \(I\setminus A_{\psi_n}\) is in \(U\), so
their intersection \(I\setminus A\) is in \(U\) by \(\sigma\)-completeness,
and \(\varnothing=A\cap(I\setminus A)\in U\), which is impossible. Existential
quantifier \(\exists w\,\psi(\bar v,w)\): if the ultraproduct satisfies it, a
witness \([g]\) gives \(A_\psi(\bar f,g)\in U\) by the induction hypothesis,
and \(A_\psi(\bar f,g)\subseteq A_{\exists w\psi}(\bar f)\). Conversely, if
\(A_{\exists w\psi}(\bar f)\in U\), choose (axiom of choice) for each \(i\) in
this set some \(g(i)\in M_i\) with \(\mathcal M_i\models\psi(\bar f(i),g(i))\),
and let \(g(i)\) be arbitrary elsewhere (structures are nonempty); then
\(A_\psi(\bar f,g)\in U\), and the induction hypothesis gives
\(\psi([\bar f],[g])\) in the ultraproduct. Universal quantifier
\(\forall w\,\psi\): if \(A_{\forall w\psi}(\bar f)\in U\) then for every
\(g\), \(A_\psi(\bar f,g)\supseteq A_{\forall w\psi}(\bar f)\) is in \(U\),
and the ultraproduct satisfies \(\psi([\bar f],[g])\) for every \([g]\). If
\(A_{\forall w\psi}(\bar f)\notin U\), choose \(g(i)\) with
\(\mathcal M_i\models\neg\psi(\bar f(i),g(i))\) for \(i\) outside that set and
arbitrary elsewhere; then \(A_\psi(\bar f,g)\subseteq A_{\forall w\psi}(\bar f)\)
is not in \(U\) (a filter is closed upward), and by the induction hypothesis
\([g]\) witnesses that \(\forall w\,\psi\) fails.
\end{proof}
\end{annexitem}

\begin{annexitem}{annex:p03-w03-p03-039-axiomatizable}{thm:p03-w03-p03-039-axiomatizable}{Classes axiomatizable in \(\mathcal L_{\omega_1,\omega}\)}
\begin{proof}
Every sentence below uses countable conjunctions and disjunctions and
finitely many quantifiers along each branch, so it lies in
\(\mathcal L_{\omega_1,\omega}\). Let \(\gamma\) be the conjunction of the
group axioms and \(\alpha\) that of the abelian group axioms.

i) \(\alpha\wedge\forall x\bigvee_{1\le n<\omega}nx=0\), where \(nx\) is the
term \(x+\cdots+x\) with \(n\) summands. An abelian group satisfies the
second conjunct exactly when every element has finite order.

ii) and iii) \(\gamma\wedge\sigma_{\mathrm{fg}}\) and
\(\gamma\wedge\neg\sigma_{\mathrm{fg}}\), with \(\sigma_{\mathrm{fg}}\) the
sentence of the annex proof of \cref{thm:p03-w03-p03-039-nadel-expressive},
item~(3), where it is shown that a group satisfies \(\sigma_{\mathrm{fg}}\)
exactly when it is finitely generated.

iv) For \(n\ge1\) let \(\theta_n(x,y)\) be
\[
\exists z_0\cdots\exists z_n\Bigl(z_0=x\wedge z_n=y\wedge\bigwedge_{k<n}\bigl(z_k<z_{k+1}\wedge\neg\exists u\,(z_k<u\wedge u<z_{k+1})\bigr)\Bigr),
\]
and let \(\sigma_{\mathbb Z}\) be the conjunction of the strict linear order
axioms, \(\forall x\exists y\,x<y\), \(\forall x\exists y\,y<x\) and
\(\forall x\forall y\bigl(x<y\to\bigvee_{1\le n<\omega}\theta_n(x,y)\bigr)\).
In \((\mathbb Z,<)\), \(\theta_n(x,y)\) holds exactly when \(y=x+n\), so
\((\mathbb Z,<)\models\sigma_{\mathbb Z}\). Conversely let
\((L,<)\models\sigma_{\mathbb Z}\). For \(x\in L\) there is \(y>x\), and
\(\theta_n(x,y)\) for some \(n\ge1\) gives \(z_1\) with \(x<z_1\) and nothing
strictly between: an immediate successor \(s(x)\), unique because \(<\) is
linear. Symmetrically, from some \(y<x\) and \(\theta_n(y,x)\) the element
\(z_{n-1}\) is an immediate predecessor \(p(x)\), and \(s(p(x))=x=p(s(x))\).
Fix \(a\in L\) and put \(h(n)=s^n(a)\) and \(h(-n)=p^n(a)\) for \(n\ge0\).
Then \(h(m)<h(m+1)=s(h(m))\) for all \(m\in\mathbb Z\), so \(h\) is strictly
increasing. If \(y>a\), then \(\theta_n(a,y)\) gives a chain
\(a=z_0<z_1<\cdots<z_n=y\) of immediate successors, so \(z_k=s^k(a)\) by
induction and \(y=h(n)\); symmetrically every \(y<a\) is some \(h(-n)\).
So \(h\) is an isomorphism \((\mathbb Z,<)\cong(L,<)\).

v) Let \(\sigma_{\mathrm{arch}}\) be the conjunction of the ordered field axioms
with \(\forall x\bigvee_{1\le n<\omega}x<n\), where \(n\) is the term
\(1+\cdots+1\) with \(n\) summands. An ordered field satisfies it exactly when
no element exceeds every natural number, which is the archimedean property.

vi) Let \(\sigma_{\mathrm{con}}\) be the conjunction of
\(\forall x\,\neg E(x,x)\), \(\forall x\forall y(E(x,y)\to E(y,x))\) and
\(\forall x\forall y\bigvee_{n<\omega}\pi_n(x,y)\), where \(\pi_n(x,y)\) is
\(\exists z_0\cdots\exists z_n\bigl(z_0=x\wedge z_n=y\wedge\bigwedge_{k<n}E(z_k,z_{k+1})\bigr)\).
A graph satisfies the last conjunct exactly when any two vertices are joined
by a path of some finite length, that is, when it is connected.

vii) The recursive sets \(p(v,w_1,\dots,w_k)\) of first-order formulas (for
all \(k\)) are countably many, since each is given by an index of a
recursive characteristic function. For each such \(p\) let
\[
\rho_p=\forall w_1\cdots\forall w_k\Bigl(\bigwedge_{p_0}\exists v\bigwedge p_0\;\to\;\exists v\bigwedge_{\varphi\in p}\varphi\Bigr),
\]
with \(p_0\) ranging over the finite subsets of \(p\), a countable set. The
sentence \(\bigwedge\mathrm{PA}\wedge\bigwedge_p\rho_p\) (with \(\mathrm{PA}\)
the countable set of Peano axioms) holds in \(\mathcal M\) exactly when
\(\mathcal M\models\mathrm{PA}\) and, for every recursive \(p\) and every
\(\bar a\), finite satisfiability of \(p(v,\bar a)\) implies its
realization, which is the definition of recursive saturation.

viii) Let \(\operatorname{Ind}(x)\) be the first-order formula
\(\exists e\,(e\in x\wedge\forall u\,\neg u\in e)\wedge\forall y\bigl(y\in x\to\exists z(z\in x\wedge\forall u(u\in z\leftrightarrow u\in y\vee u=y))\bigr)\)
and \(\operatorname{Om}(x)\) the formula
\(\operatorname{Ind}(x)\wedge\forall y(\operatorname{Ind}(y)\to\forall u(u\in x\to u\in y))\);
in a model of \(\mathrm{ZFC}\) exactly one element satisfies
\(\operatorname{Om}\), namely \(\omega^{\mathcal M}\). Let
\(\delta_0(y)=\forall u\,\neg u\in y\) and
\(\delta_{n+1}(y)=\exists z\bigl(\delta_n(z)\wedge\forall u(u\in y\leftrightarrow u\in z\vee u=z)\bigr)\).
By induction on \(n\), using extensionality, in a model of \(\mathrm{ZFC}\)
the element \(n^{\mathcal M}\) is the only element satisfying \(\delta_n\).
Let \(\sigma=\bigwedge\mathrm{ZFC}\wedge\forall x\forall y\bigl(\operatorname{Om}(x)\wedge y\in x\to\bigvee_{n<\omega}\delta_n(y)\bigr)\),
a countable conjunction. A model of \(\mathrm{ZFC}\) satisfies the last
conjunct exactly when every element of \(\omega^{\mathcal M}\) is some
\(n^{\mathcal M}\), which is the definition of an \(\omega\)-model.

ix) Since the vocabulary is countable, \(T\) and \(p(\bar v)\) are countable.
The set \(T\cup\{\sigma_p\}\), with \(\sigma_p\) the sentence of the annex
proof of \cref{thm:p03-w03-p03-039-nadel-expressive}, item~(5), for
\(\Phi=p\), axiomatizes the models of \(T\) omitting \(p\); for countable
\(T\) the single sentence \(\bigwedge T\wedge\sigma_p\) does.
\end{proof}
\end{annexitem}

\begin{annexitem}{annex:p03-w03-p03-039-upward-ls-failure}{thm:p03-w03-p03-039-upward-ls-failure}{The upward L\"owenheim--Skolem theorem fails}
\begin{proof}
(1) By item iv) of \cref{thm:p03-w03-p03-039-axiomatizable} the models of
\(\sigma_{\mathbb Z}\) are the orders isomorphic to \((\mathbb Z,<)\), which
is countably infinite.

(2) If \(a_1,\dots,a_n\) generate \(G\), every element is \(t(a_1,\dots,a_n)\)
for a term \(t\) in \(n\) variables (annex proof of
\cref{thm:p03-w03-p03-039-nadel-expressive}, item~(3)), and there are
countably many terms. \((\mathbb Z,+)\) is generated by \(1\) and is infinite.

(3) Let \(F\) be an archimedean ordered field. Its prime field is an ordered
copy of \(\mathbb Q\), since \(n\cdot1>0\) for \(n\ge1\). For \(x\in F\) let
\(C(x)=\{q\in\mathbb Q:q<x\}\). By the archimedean property there are
naturals \(k,l\) with \(-x<k\) and \(x<l\), so \(-k\in C(x)\) and \(C(x)\) is
bounded above by \(l\); let \(h(x)=\sup C(x)\in\mathbb R\). Let \(x<y\). Choose
\(n\ge1\) with \(1/(y-x)<n\), so \(1/n<y-x\). The set of integers \(m\) with
\(nx<m\) is nonempty and bounded below (archimedean property applied to
\(nx\) and \(-nx\)); let \(m\) be its least element. Then \(m-1\le nx<m\), so
\(x<m/n\le x+1/n<y\). Applying the same argument to \(m/n<y\) gives a
rational \(q'\) with \(m/n<q'<y\). Then \(C(x)\subseteq\{q:q<m/n\}\), so
\(h(x)\le m/n\), while \(q'\in C(y)\), so \(h(y)\ge q'>m/n\). Thus \(h\) is a
strictly increasing map \(F\to\mathbb R\), in particular injective, and
\(\lvert F\rvert\le\lvert\mathbb R\rvert=2^{\aleph_0}\). \(\mathbb Q\) and
\(\mathbb R\) satisfy the archimedean property.
\end{proof}
\end{annexitem}

\begin{annexitem}{annex:p03-w03-p03-039-formula-count}{thm:p03-w03-p03-039-formula-count}{Counting formulas of \(\mathcal L_{\kappa,\omega}\)}
\begin{proof}
Let \(S_0\) be the set of atomic formulas: terms are finite strings over the
symbols of \(\tau\) and the variables \(v_\xi\) (\(\xi<\kappa\)), so there are
at most \(\kappa\) terms and at most \(\kappa\) atomic formulas. Let
\(S_{\xi+1}\) be \(S_\xi\) together with \(\neg\varphi\), \(\exists v_\zeta\varphi\),
\(\forall v_\zeta\varphi\) for \(\varphi\in S_\xi\) and \(\zeta<\kappa\), and
\(\bigwedge X\), \(\bigvee X\) for \(X\subseteq S_\xi\) with \(\lvert X\rvert<\kappa\);
and \(S_\lambda=\bigcup_{\xi<\lambda}S_\xi\) at limits. By induction
\(\lvert S_\xi\rvert\le2^\kappa\): at successors there are at most
\(\lvert S_\xi\rvert^{\kappa}\le(2^\kappa)^\kappa=2^\kappa\) subsets of size
less than \(\kappa\), and at limits \(\xi\le\kappa\) a union of at most
\(\kappa\) sets of size at most \(2^\kappa\) has size at most \(2^\kappa\).
Every member of \(S_\kappa\) is a formula of \(\mathcal L_{\kappa,\omega}(\tau)\).
Conversely \(S_\kappa\) contains the atomic formulas and is closed under the
formation rules: the unary rules are clear, and if \(X\subseteq S_\kappa\)
with \(\lvert X\rvert<\kappa\), each member lies in some \(S_\xi\) with
\(\xi<\kappa\), the supremum \(\eta\) of these fewer than \(\kappa\) ordinals
is below \(\kappa\) because \(\kappa\) is regular, so \(X\subseteq S_\eta\)
and \(\bigwedge X,\bigvee X\in S_{\eta+1}\). Since the formulas form the least
class with these closure properties, they all lie in \(S_\kappa\), which has
cardinality at most \(2^\kappa\).
\end{proof}
\end{annexitem}

\begin{annexitem}{annex:p03-w03-p03-039-ordinal-equivalence}{thm:p03-w03-p03-039-ordinal-equivalence}{Ordinals that \(\mathcal L_{\kappa,\omega}\) cannot tell apart}
\begin{proof}
Let \(F\) be the set of sentences of \(\mathcal L_{\kappa,\omega}(\{<\})\);
by \cref{thm:p03-w03-p03-039-formula-count}, \(\lvert F\rvert\le2^\kappa\).
The map sending a nonzero ordinal \(\alpha\) to
\(\{\varphi\in F:(\alpha,<)\models\varphi\}\) takes values in the power set
of \(F\), which has at most \(2^{2^\kappa}\) members. There are
\((2^{2^\kappa})^+\) nonzero ordinals below \((2^{2^\kappa})^+\), so the map
is not injective on them: some \(\alpha\ne\beta\) have the same
\(\mathcal L_{\kappa,\omega}\)-theory, that is,
\((\alpha,<)\equiv_{\kappa,\omega}(\beta,<)\). By
\cref{thm:p03-w03-p03-039-nadel-ordinals}, \((\alpha,<)\models\sigma_\alpha\)
and \((\beta,<)\not\models\sigma_\alpha\) (the latter would give
\((\beta,<)\cong(\alpha,<)\), and distinct ordinals are not isomorphic), so
\((\alpha,<)\not\equiv_{\infty,\omega}(\beta,<)\).
\end{proof}
\end{annexitem}

\begin{annexitem}{annex:p03-w03-p03-039-pc-examples}{thm:p03-w03-p03-039-pc-examples}{Six \(\mathrm{PC}_{\omega_1,\omega}\)-classes}
\begin{proof}
In each case a vocabulary \(\tau^*\) and a sentence \(\varphi\) are given, and
the class is shown to be the class of \(\tau\)-reducts of models of
\(\varphi\). Let \(\gamma\) be the conjunction of the group axioms.

a) \(\tau^*\) adds \(<\); \(\varphi\) is \(\gamma\), the strict linear order
axioms and \(\forall x\forall y\forall z\bigl(x<y\to zx<zy\wedge xz<yz\bigr)\).
A model is exactly a group with an order of the required kind. First-order.

b) \(\tau^*\) adds a unary \(B\). For \(n\ge0\) let \(W_n\) be the set of
reduced words in the letters \(y_1^{\pm1},\dots,y_n^{\pm1}\), read as group
terms (the empty word is \(1\), and \(W_0\) contains only it); \(W_n\) is
countable. Let \(\delta_n(\bar y)\) say
\(\bigwedge_{i}B(y_i)\wedge\bigwedge_{i<j}y_i\ne y_j\) (\(\delta_0\) is the
empty conjunction, and \(\exists\bar y\) is then empty). Let \(\varphi\) be
the conjunction of \(\gamma\),
\(\forall x\bigvee_{n\ge0}\bigvee_{w\in W_n}\exists\bar y(\delta_n(\bar y)\wedge x=w(\bar y))\)
and, for all \(n\ge0\) and distinct \(u,w\in W_n\),
\(\forall\bar y(\delta_n(\bar y)\to u(\bar y)\ne w(\bar y))\). A reduced word
over \(B\) uses finitely many distinct elements \(b_1,\dots,b_n\) of \(B\) and
is \(w(\bar b)\) for exactly one \(w\in W_n\) once the \(b_i\) are listed;
two reduced words over \(B\) can be listed over one common injective tuple,
and then they are distinct as words over \(B\) exactly when the
corresponding members of \(W_n\) are distinct. So the second conjunct says
that every element is the value of a reduced word over \(B\), and the third
that distinct reduced words over \(B\) have distinct values: \(G\) is free on
\(B\). Conversely a group free on \(B\) satisfies \(\varphi\) with
\(B\) interpreted as the basis.

c) \(\tau^*\) adds a ternary function symbol \(F\); \(\varphi\) is the strict
linear order axioms together with
\(\forall a\forall b\bigl(F(a,b,a)=b\wedge\forall x\forall y(x<y\to F(a,b,x)<F(a,b,y))\wedge\forall y\exists x\,F(a,b,x)=y\bigr)\).
In a model, \(x\mapsto F(a,b,x)\) is strictly increasing (hence injective
and order-reflecting, the order being linear) and surjective, an automorphism
taking \(a\) to \(b\). Conversely, in a \(1\)-transitive order choose (axiom
of choice) an automorphism \(\pi_{a,b}\) with \(\pi_{a,b}(a)=b\) for each pair
and put \(F(a,b,x)=\pi_{a,b}(x)\). First-order.

d) \(\tau^*\) adds a unary \(B\); \(\varphi\) is the conjunction of the field
axioms, \(\bigwedge_{n\ge1}n\cdot1\ne0\), for every \(n\ge1\) and every
nonzero \(f\in\mathbb Z[y_1,\dots,y_n]\) the sentence
\(\forall\bar y(\delta_n(\bar y)\to f(\bar y)\ne0)\), and
\(\forall x\bigvee_{n\ge0}\bigvee_{f,g\in\mathbb Z[y_1,\dots,y_n]}\exists\bar y\bigl(\delta_n(\bar y)\wedge g(\bar y)\ne0\wedge x\cdot g(\bar y)=f(\bar y)\bigr)\),
where \(\delta_n\) is as in b), \(\delta_0\) is the empty conjunction and
\(\mathbb Z[y_1,\dots,y_0]=\mathbb Z\). In a model the field has characteristic
\(0\), so its prime field is \(\mathbb Q\); a nonzero polynomial over
\(\mathbb Q\) becomes one over \(\mathbb Z\) after clearing denominators, so
\(B\) is algebraically independent over \(\mathbb Q\); and every element is a
quotient of integer polynomials in finitely many members of \(B\), so the
field is \(\mathbb Q(B)\). Conversely \(\mathbb Q(B)\) with \(B\)
algebraically independent satisfies \(\varphi\), every element of
\(\mathbb Q(B)\) being such a quotient. The case \(B=\varnothing\), \(K=\mathbb Q\), is the disjunct \(n=0\).

e) \(\tau^*\) adds a unary \(S\); \(\varphi\) is the ordered field axioms,
\(\exists x\,S(x)\), \(\exists b\forall x(S(x)\to x\le b)\) and
\(\forall b\bigl(\forall x(S(x)\to x\le b)\to\exists b'(b'<b\wedge\forall x(S(x)\to x\le b'))\bigr)\).
The last three say that \(S\) is nonempty, bounded above, and that no upper
bound is least. First-order.

f) \(\tau^*\) adds a unary \(Z\); \(\varphi\) is the ordered field axioms,
\(Z(0)\wedge Z(1)\), closure of \(Z\) under \(+\), \(-\) and \(\cdot\),
\(\neg\exists z(Z(z)\wedge0<z\wedge z<1)\) and
\(\forall x\exists z\bigl(Z(z)\wedge-1\le x-z\wedge x-z\le1\bigr)\). These
are the clauses of the definition of an integer part. First-order.
\end{proof}
\end{annexitem}

\begin{annexitem}{annex:p03-w03-p03-039-silver}{thm:p03-w03-p03-039-silver}{Silver's class}
\begin{proof}
(1) Let \(\tau^*=\{U,E,f\}\) with \(E\) binary and \(f\) unary, and let
\(\varphi\) be the conjunction of
\(\forall x\forall y\bigl(\forall u(U(u)\to(E(u,x)\leftrightarrow E(u,y)))\to x=y\bigr)\),
\(\forall x\,\neg U(f(x))\) and \(\forall x\forall y(f(x)=f(y)\to x=y)\). In a
model, \(x\mapsto\{u\in U:E(u,x)\}\) is an injection of \(M\) into the power
set of \(U(\mathcal M)\), so \(\lvert M\rvert\le2^{\lvert U(\mathcal M)\rvert}\),
and \(f\) injects \(M\) into \(M\setminus U(\mathcal M)\), so
\(\lvert M\rvert\le\lvert M\setminus U(\mathcal M)\rvert\le\lvert M\rvert\).
Conversely, for \(\mathcal M\in\mathcal K\) choose an injection
\(g:M\to\mathcal P(U(\mathcal M))\) and put \(E(u,x)\) iff \(u\in g(x)\), and
choose an injection \(f:M\to M\setminus U(\mathcal M)\), which exists because
the two sets have the same cardinality.

(2) A \(\{U\}\)-structure is determined up to isomorphism by the pair
\((\lvert U(\mathcal M)\rvert,\lvert M\setminus U(\mathcal M)\rvert)\): a
bijection on each part combines to an isomorphism. A structure of
cardinality \(\kappa\) is in \(\mathcal K\) exactly when
\(\lvert M\setminus U\rvert=\kappa\) and \(\kappa\le2^{\lambda}\), where
\(\lambda=\lvert U\rvert\le\kappa\). So the models of cardinality \(\kappa\)
correspond to the cardinals \(\lambda\le\kappa\) with \(\kappa\le2^\lambda\);
\(\lambda=\kappa\) is always one of them. Hence \(\mathcal K\) has exactly one
model of cardinality \(\kappa\) iff \(2^\lambda<\kappa\) for every
\(\lambda<\kappa\), that is, iff \(\kappa\) is a strong limit.

It remains to identify the strong limits. \(\aleph_0=\beth_0\) is one, since
\(2^n\) is finite. If \(\delta\) is a limit ordinal and
\(\lambda<\beth_\delta=\sup_{\beta<\delta}\beth_\beta\), then
\(\lambda<\beth_\beta\) for some \(\beta<\delta\) and
\(2^\lambda\le2^{\beth_\beta}=\beth_{\beta+1}<\beth_\delta\). Conversely let
\(\kappa\) be an infinite strong limit and \(\alpha\) least with
\(\kappa\le\beth_\alpha\). If \(\alpha=\beta+1\), then \(\beth_\beta<\kappa\)
and \(\beth_{\beta+1}=2^{\beth_\beta}<\kappa\), a contradiction. If
\(\alpha=0\), then \(\kappa=\aleph_0\). If \(\alpha\) is a limit, then
\(\beth_\beta<\kappa\) for all \(\beta<\alpha\), so
\(\beth_\alpha=\sup_{\beta<\alpha}\beth_\beta\le\kappa\le\beth_\alpha\).
\end{proof}
\end{annexitem}

\begin{annexitem}{annex:p03-w03-p03-039-large-models}{thm:p03-w03-p03-039-large-models}{A countable vocabulary with only large models}
\begin{proof}
Each \(\theta_S\) is a sentence of \(\mathcal L_{\omega_1,\omega}\). The
structure with universe \(\mathcal P(\omega)\) and
\(U_n=\{S:n\in S\}\) satisfies \(\theta_S\) with witness \(S\). If
\(\mathcal M\models T\), choose (axiom of choice) a witness \(x_S\) for each
\(\theta_S\). For \(S\ne S'\) pick \(n\) in their symmetric difference; then
exactly one of \(x_S,x_{S'}\) is in \(U_n\), so \(x_S\ne x_{S'}\). Thus
\(S\mapsto x_S\) is injective and \(\lvert M\rvert\ge2^{\aleph_0}\).
\end{proof}
\end{annexitem}

\begin{annexitem}{annex:p03-w03-p03-039-definable-approximants}{thm:p03-w03-p03-039-definable-approximants}{The approximants and the Scott rank are definable}
\begin{proof}
(1) By recursion on \(\beta<\omega_1\), simultaneously for all \(n\), with
\(x,y\) variables not among \(\bar v,\bar w\): \(S^n_0(\bar v,\bar w)\) is
\(\bigwedge_\varphi\bigl(\varphi(\bar v)\leftrightarrow\varphi(\bar w)\bigr)\),
over the atomic \(\tau\)-formulas \(\varphi(x_1,\dots,x_n)\), a countable set
because \(\tau\) is countable;
\(S^n_{\beta+1}=\forall x\exists y\,S^{n+1}_\beta(\bar vx,\bar wy)\wedge\forall y\exists x\,S^{n+1}_\beta(\bar vx,\bar wy)\);
and \(S^n_\beta=\bigwedge_{\gamma<\beta}S^n_\gamma\) for a limit \(\beta\), a
countable conjunction because \(\beta<\omega_1\). Each \(S^n_\beta\) is in
\(\mathcal L_{\omega_1,\omega}(\tau)\). The claim follows by induction on
\(\beta\), since the three clauses restate the three clauses of
\cref{def:p03-w03-p03-039-approximants} for \(\mathcal N=\mathcal M\).

(2) Let \(\Gamma_\alpha\) be the set of pairs of equal-length tuples from
\(M\) related by \(\sim_\alpha\), and \(\theta\) the cardinality of the set of
all such pairs (infinite, as there are tuples of every length). By
\cref{thm:p03-w03-p03-039-approximant-basics} the \(\Gamma_\alpha\) decrease.
If \(\Gamma_\alpha\ne\Gamma_{\alpha+1}\) for every \(\alpha<\theta^+\), the
choice of a pair in \(\Gamma_\alpha\setminus\Gamma_{\alpha+1}\) for each
\(\alpha\) is injective (a pair chosen at \(\beta>\alpha\) lies in
\(\Gamma_{\alpha+1}\)), which is impossible. At an \(\alpha\) with
\(\Gamma_\alpha=\Gamma_{\alpha+1}\), every \(\bar a\) satisfies the condition
defining \(r(\bar a)\); so the least ordinal satisfying it exists. Next,
\(r(\bar a)\ge\gamma\) holds iff the condition fails at every \(\delta<\gamma\),
that is, iff for every \(\delta<\gamma\) some \(\bar b\) has
\((\mathcal M,\bar a)\sim_\delta(\mathcal M,\bar b)\) and not
\((\mathcal M,\bar a)\sim_{\delta+1}(\mathcal M,\bar b)\). And
\(\operatorname{SR}(\mathcal M)=\sup_{\bar a}(r(\bar a)+1)\ge\alpha\) iff for
every \(\gamma<\alpha\) some \(\bar a\) has \(r(\bar a)+1>\gamma\), that is,
\(r(\bar a)\ge\gamma\). With \(\bar v,\bar w\) of length \(n\), let
\[
\sigma_\alpha=\bigwedge_{\gamma<\alpha}\ \bigvee_{n<\omega}\ \exists\bar v\bigwedge_{\delta<\gamma}\exists\bar w\bigl(S^n_\delta(\bar v,\bar w)\wedge\neg S^n_{\delta+1}(\bar v,\bar w)\bigr).
\]
All conjunctions and disjunctions are countable because \(\alpha<\omega_1\),
and by item~(1) and the two equivalences just shown,
\(\mathcal M\models\sigma_\alpha\) iff \(\operatorname{SR}(\mathcal M)\ge\alpha\).
\end{proof}
\end{annexitem}

\begin{annexitem}{annex:p03-w03-p03-039-countable-substructures}{thm:p03-w03-p03-039-countable-substructures}{Countably generated substructures transfer}
\begin{proof}
By \cref{thm:p03-w03-p03-039-karp} there is a partial isomorphism system
\(P\) between \(\mathcal M\) and \(\mathcal N\). Let \(\mathcal M_0\) be
generated by \(\{a_0,a_1,\ldots\}\) (repeat elements if the set is finite).
Choose \(f_0\in P\) and, by clause~1 of
\cref{def:p03-w03-p03-039-partial-iso}, \(f_{k+1}\in P\) extending \(f_k\)
with \(a_k\in\operatorname{dom}(f_{k+1})\). The union \(f\) of the chain is a
function defined on \(D=\{a_k:k\in\omega\}\), and it preserves and reflects
atomic formulas on tuples from its domain, since such a tuple lies in some
\(\operatorname{dom}(f_k)\). Every element of \(M_0\) is \(t^{\mathcal M}(\bar a)\)
for a term \(t\) and a tuple \(\bar a\) from \(D\). Put
\(h(t^{\mathcal M}(\bar a))=t^{\mathcal N}(f(\bar a))\). Since
\(t(\bar a)=s(\bar a')\) is an atomic statement about a tuple from \(D\), it
holds in \(\mathcal M\) iff the corresponding statement holds in
\(\mathcal N\) at \(f(\bar a),f(\bar a')\); so \(h\) is well defined and
injective. It commutes with the function symbols and constants by its
definition, and it preserves and reflects each relation \(R\), since
\(R(t_1(\bar a),\dots)\) is atomic. So \(h\) is an embedding of
\(\mathcal M_0\) into \(\mathcal N\), and its image is a substructure of
\(\mathcal N\) isomorphic to \(\mathcal M_0\).
\end{proof}
\end{annexitem}

\begin{annexitem}{annex:p03-w03-p03-039-subgroups-free}{thm:p03-w03-p03-039-subgroups-free}{Countable subgroups of free abelian groups are free}
\begin{proof}
Let \(F\) be free abelian with basis \(B\) and \(S\le F\) countable. Each
element of \(S\) is an integer combination of finitely many members of \(B\);
let \(B_0\) be the countable set of members of \(B\) so used, enumerated as
\(e_0,e_1,\ldots\) (a finite or infinite list). Let \(F_n\) be the subgroup
generated by \(e_0,\dots,e_{n-1}\) and \(S_n=S\cap F_n\); then
\(S=\bigcup_nS_n\). The map \(\pi_n:F_{n+1}\to\mathbb Z\) taking the
coefficient of \(e_n\) is a homomorphism with kernel \(F_n\), and
\(\pi_n(S_{n+1})=m_n\mathbb Z\) for some \(m_n\ge0\). If \(m_n\ne0\) choose
\(s_n\in S_{n+1}\) with \(\pi_n(s_n)=m_n\). For \(x\in S_{n+1}\),
\(\pi_n(x)=km_n\) for some \(k\), and \(x-ks_n\in S_{n+1}\cap\ker\pi_n=S_n\);
and \(ks_n\in S_n\) forces \(km_n=0\), so \(k=0\). Hence \(S_{n+1}=S_n\oplus\mathbb Zs_n\)
when \(m_n\ne0\), and \(S_{n+1}=S_n\) when \(m_n=0\). By induction on \(n\),
\(C_n=\{s_k:k<n,\ m_k\ne0\}\) is a basis of \(S_n\) (\(S_0=\{0\}\) has the
empty basis). Then \(C=\bigcup_nC_n\) generates \(S\), and a nontrivial
integer relation among members of \(C\) would involve finitely many, all in
some \(C_n\), contradicting that \(C_n\) is a basis. So \(S\) is free.

If \(G\) is free abelian and \(H\equiv_{\infty,\omega}G\), then \(H\) is an
abelian group, the abelian group axioms being sentences true in \(G\). A
countable subgroup \(H_0\) of \(H\) is a countably generated substructure in
the vocabulary \(\{+,-,0\}\), so by
\cref{thm:p03-w03-p03-039-countable-substructures} (with the roles of the
structures exchanged) it is isomorphic to a countable subgroup of \(G\),
which is free by the first part.
\end{proof}
\end{annexitem}

\begin{annexitem}{annex:p03-w03-p03-039-aleph-one-free}{thm:p03-w03-p03-039-aleph-one-free}{\(\aleph_1\)-free groups and the free group of countable rank}
\begin{proof}
For \(k\in\mathbb Z\) and a variable \(x\) let \(kx\) be the term
\(x+\cdots+x\) (\(k\) summands) for \(k>0\), \(0\) for \(k=0\), and
\(-(\lvert k\rvert x)\) for \(k<0\). For \(n\ge1\) let \(\psi_n\) be
\[
\exists x_1\cdots\exists x_n\Bigl(\forall y\bigvee_{\bar k\in\mathbb Z^n}y=k_1x_1+\cdots+k_nx_n\;\wedge\bigwedge_{\bar k\in\mathbb Z^n\setminus\{\bar0\}}k_1x_1+\cdots+k_nx_n\ne0\Bigr),
\]
and \(\psi_0=\forall y\,y=0\). In an abelian group the first conjunct says
that \(x_1,\dots,x_n\) generate and the second that they are independent
(which forces them to be distinct and nonzero); so an abelian group
satisfies \(\psi_n\) iff it is free on a set of \(n\) elements. Let
\(\chi=\neg\bigvee_{n<\omega}\psi_n\); an abelian group satisfies \(\chi\) iff
it is not free on a finite set.

\(G\) is countably infinite, so its Scott sentence \(\Phi^G\) is in
\(\mathcal L_{\omega_1,\omega}(\{+,-,0\})\)
(\cref{def:p03-w03-p03-039-scott-sentence}). Let \(A\) be the least countable
fragment containing \(\Phi^G\), \(\chi\) and the abelian group axioms
(\cref{thm:p03-w03-p03-039-countable-fragment-size}); its formulas have
finitely many free variables. By \cref{thm:p03-w03-p03-039-downward-ls} there
is a countable \(\mathcal H_1\prec_A\mathcal H\); it is a substructure, hence
a countable subgroup, hence free with some basis \(C\) because \(H\) is
\(\aleph_1\)-free. Since \(H\models\chi\) and \(\chi\in A\),
\(\mathcal H_1\models\chi\), so \(C\) is infinite; it is countable, so there
is a bijection between \(C\) and a basis of \(G\), which extends to an
isomorphism \(\mathcal H_1\cong G\). Hence \(\mathcal H_1\models\Phi^G\)
(\cref{thm:p03-w03-p03-039-iso-invariance} and \(G\models\Phi^G\)), so
\(H\models\Phi^G\) because \(\Phi^G\in A\), and
\(H\equiv_{\infty,\omega}G\) by \cref{thm:p03-w03-p03-039-scott}.
\end{proof}
\end{annexitem}

\begin{annexitem}{annex:p03-w03-p03-039-aleph-one-free-axiom}{thm:p03-w03-p03-039-aleph-one-free-axiom}{\(\aleph_1\)-freeness is axiomatizable}
\begin{proof}
With \(G\), \(\Phi^G\) and \(\psi_n\) as in the proof of
\cref{thm:p03-w03-p03-039-aleph-one-free}, let the sentence be the
conjunction of the abelian group axioms with
\(\Phi^G\vee\bigvee_{n<\omega}\psi_n\). If \(A\) is \(\aleph_1\)-free and free
on a finite set of \(n\) elements, \(A\models\psi_n\); otherwise
\(A\equiv_{\infty,\omega}G\) by \cref{thm:p03-w03-p03-039-aleph-one-free}, so
\(A\models\Phi^G\). Conversely, if \(A\models\Phi^G\) then
\(A\equiv_{\infty,\omega}G\) by \cref{thm:p03-w03-p03-039-scott} and \(A\) is
\(\aleph_1\)-free by \cref{thm:p03-w03-p03-039-subgroups-free}; if
\(A\models\psi_n\), \(A\) is free and its countable subgroups are free by
the same lemma.
\end{proof}
\end{annexitem}

\begin{annexitem}{annex:p03-w03-p03-039-countable-fragment-size}{thm:p03-w03-p03-039-countable-fragment-size}{Fragments of a given size}
\begin{proof}
For an atomic formula the set of subformulas and the set of variables are
finite. A formula of \(\mathcal L_{\kappa^+,\omega}\) has at most \(\kappa\)
subformulas by \cref{thm:p03-w03-p03-039-subformula-count} (\(\kappa^+\) is
regular), and hence at most \(\kappa\) variables, each occurring in an atomic
subformula. Let \(S_0\) be \(T\) together with the atomic formulas of \(\tau\)
with variables in \(V\); there are at most \(\kappa\) terms over \(\tau\) and
\(V\), so \(\lvert S_0\rvert\le\kappa\). Given \(S_k\) of cardinality at most
\(\kappa\) whose members lie in \(\mathcal L_{\kappa^+,\omega}\), have at most
\(\kappa\) subformulas each, use only variables from \(V\) and have finitely
many free variables, let \(S_{k+1}\) add: the subformulas of members of
\(S_k\); the results of substituting a term over \(\tau\) and \(V\) for a free
variable of a member of \(S_k\); the formal negations of members of
\(S_k\); and \(\neg\varphi\), \(\varphi\wedge\psi\), \(\varphi\vee\psi\),
\(\exists v\,\varphi\), \(\forall v\,\varphi\) for \(\varphi,\psi\in S_k\) and
\(v\in V\). Each addition is indexed by at most \(\kappa\cdot\kappa=\kappa\)
data, so \(\lvert S_{k+1}\rvert\le\kappa\). The four properties carry over:
a subformula of a formula with finitely many free variables has finitely
many (\cref{thm:p03-w03-p03-039-nadel-free-variables}); substitution replaces
one free variable by the finitely many variables of a term and maps the
subformulas of \(\varphi\) onto those of the result; formal negation keeps
the free variables (\cref{thm:p03-w03-p03-039-formal-negation}) and its
subformulas are formal negations of subformulas, together with subformulas,
of \(\varphi\); and the finitary steps add one formula on top of at most two.
Then \(A=\bigcup_kS_k\) has cardinality at most \(\kappa\), and it is a
fragment with witness \(V\) containing \(T\), by the closing paragraph of the
proof of \cref{thm:p03-w03-p03-039-countable-fragment}, which uses nothing
about cardinality. For \(\kappa=\aleph_0\) and \(T\) a set of sentences this
is the fragment of \cref{thm:p03-w03-p03-039-countable-fragment}, and its
formulas have finitely many free variables.
\end{proof}
\end{annexitem}

\begin{annexitem}{annex:p03-w03-p03-039-finite-structures}{thm:p03-w03-p03-039-finite-structures}{Finite structures and countable models}
\begin{proof}
(1) List \(M=\{a_1,\dots,a_n\}\) without repetition and let \(D(x_1,\dots,x_n)\)
be the set of atomic and negated atomic \(\tau\)-formulas in
\(x_1,\dots,x_n\) true of \(\bar a\); it is countable since \(\tau\) is. Let
\(\chi_{\mathcal M}=\exists x_1\cdots\exists x_n\bigl(\bigwedge_{i<j}x_i\ne x_j\wedge\forall y\bigvee_iy=x_i\wedge\bigwedge D\bigr)\).
\(\mathcal M\) satisfies it with witnesses \(\bar a\), and so does every
isomorphic copy (\cref{thm:p03-w03-p03-039-iso-invariance}). If
\(\mathcal N\models\chi_{\mathcal M}\) with witnesses \(\bar b\), then
\(N=\{b_1,\dots,b_n\}\) without repetition, and \(h(a_i)=b_i\) is a
bijection. For a constant \(c\) with \(c^{\mathcal M}=a_k\), the formula
\(c=x_k\) is in \(D\), so \(c^{\mathcal N}=b_k\); for a function symbol
\(F\) with \(F(a_{i_1},\dots)=a_k\), the formula \(F(x_{i_1},\dots)=x_k\) is
in \(D\); and each relation holds or fails at \(\bar a\) according to \(D\).
So \(h\) is an isomorphism.

(2) Let \(\mathcal M\models\sigma\) and let \(A\) be the least countable
fragment containing \(\sigma\); its formulas have finitely many free
variables (\cref{thm:p03-w03-p03-039-countable-fragment-size}). By
\cref{thm:p03-w03-p03-039-downward-ls} with \(X=\varnothing\) there is
\(\mathcal N\prec_A\mathcal M\) with \(\lvert N\rvert\le\aleph_0\), and
\(\mathcal N\models\sigma\).
\end{proof}
\end{annexitem}

\begin{annexitem}{annex:p03-w03-p03-039-complete-extension}{thm:p03-w03-p03-039-complete-extension}{Every satisfiable sentence has a complete strengthening}
\begin{proof}
By \cref{thm:p03-w03-p03-039-finite-structures}~(2), \(\psi\) has a model
\(\mathcal M_0\) with \(\lvert M_0\rvert\le\aleph_0\). If \(M_0\) is finite,
let \(\varphi=\chi_{\mathcal M_0}\): its models are the copies of
\(\mathcal M_0\), so they all satisfy \(\psi\) and they agree on every
sentence (\cref{thm:p03-w03-p03-039-iso-invariance}); \(\varphi\) is
satisfiable, hence complete. If \(M_0\) is infinite, let
\(\varphi=\Phi^{\mathcal M_0}\), which is in \(\mathcal L_{\omega_1,\omega}(\tau)\)
because \(\tau\) and \(M_0\) are countable, and is complete by
\cref{thm:p03-w03-p03-039-scott-complete}; every model of it is
\(\equiv_{\infty,\omega}\mathcal M_0\) by \cref{thm:p03-w03-p03-039-scott},
hence satisfies \(\psi\).
\end{proof}
\end{annexitem}

\begin{annexitem}{annex:p03-w03-p03-039-complete-categorical}{thm:p03-w03-p03-039-complete-categorical}{Complete means \(\aleph_0\)-categorical}
\begin{proof}
Let \(\varphi\) be complete. It has a model of cardinality at most
\(\aleph_0\) by \cref{thm:p03-w03-p03-039-finite-structures}~(2). If
\(\mathcal M,\mathcal N\) are two such models, then for every sentence
\(\theta\), \(\varphi\models\theta\) or \(\varphi\models\neg\theta\), so
\(\mathcal M\) and \(\mathcal N\) agree on \(\theta\); thus
\(\mathcal M\equiv_{\infty,\omega}\mathcal N\), \(\mathcal M\cong_p\mathcal N\)
by \cref{thm:p03-w03-p03-039-karp}, and
\(\mathcal M\cong\mathcal N\) by \cref{thm:p03-w03-p03-039-countable-back-forth}.

Conversely let \(\mathcal M_0\) be the unique model of \(\varphi\) of
cardinality at most \(\aleph_0\), and let \(\theta_0\) be
\(\chi_{\mathcal M_0}\) if \(M_0\) is finite and \(\Phi^{\mathcal M_0}\)
otherwise; in both cases \(\theta_0\in\mathcal L_{\omega_1,\omega}(\tau)\),
\(\mathcal M_0\models\theta_0\), and \(\theta_0\) is complete (proof of
\cref{thm:p03-w03-p03-039-complete-extension}). If \(\varphi\wedge\neg\theta_0\)
had a model, it would have one of cardinality at most \(\aleph_0\), which
would be isomorphic to \(\mathcal M_0\) and so satisfy \(\theta_0\). Hence
\(\varphi\models\theta_0\), and for every sentence \(\theta\),
\(\theta_0\models\theta\) or \(\theta_0\models\neg\theta\) gives the same for
\(\varphi\). As \(\varphi\) is satisfiable, it is complete.
\end{proof}
\end{annexitem}

\begin{annexitem}{annex:p03-w03-p03-039-categorical-equivalence}{thm:p03-w03-p03-039-categorical-equivalence}{Above a categoricity cardinal all models are \(\equiv_{\infty,\omega}\)}
\begin{proof}
Let \(\mathcal N_\kappa\) be the model of \(\psi\) of cardinality \(\kappa\),
and \(\Phi=\Phi^{\mathcal N_\kappa}\), a sentence of
\(\mathcal L_{\kappa^+,\omega}(\tau)\) (\cref{def:p03-w03-p03-039-scott-sentence}).
It suffices to show that every \(\mathcal M\models\psi\) with
\(\lvert M\rvert\ge\kappa\) satisfies \(\Phi\): then
\(\mathcal M\equiv_{\infty,\omega}\mathcal N_\kappa\) by
\cref{thm:p03-w03-p03-039-scott}, and two such models are
\(\equiv_{\infty,\omega}\)-equivalent to each other.

Let \(V\) be a set of \(\kappa\) variables containing those of \(\psi\) and
\(\Phi\) (at most \(\kappa\) of them, by
\cref{thm:p03-w03-p03-039-countable-fragment-size}). By
\cref{thm:p03-w03-p03-039-countable-fragment-size} there is a fragment \(A\)
with witness \(V\), \(\psi,\Phi\in A\), \(\lvert A\rvert\le\kappa\), all of
whose formulas have finitely many free variables. Choose \(X\subseteq M\)
with \(\lvert X\rvert=\kappa\). By \cref{thm:p03-w03-p03-039-downward-ls}
there is \(\mathcal M'\prec_A\mathcal M\) with \(X\subseteq M'\) and
\(\lvert M'\rvert\le\kappa\), so \(\lvert M'\rvert=\kappa\). Since
\(\psi\in A\), \(\mathcal M'\models\psi\), so
\(\mathcal M'\cong\mathcal N_\kappa\) and \(\mathcal M'\models\Phi\)
(\(\mathcal N_\kappa\models\Phi\) by the proof of
\cref{thm:p03-w03-p03-039-scott}, and
\cref{thm:p03-w03-p03-039-iso-invariance}). Since \(\Phi\in A\),
\(\mathcal M\models\Phi\).
\end{proof}
\end{annexitem}

\begin{annexitem}{annex:p03-w03-p03-039-elementary-chain}{thm:p03-w03-p03-039-elementary-chain}{Unions of \(A\)-elementary chains}
\begin{proof}
Each \(\mathcal M_i\) is a substructure of each later \(\mathcal M_j\), so the
union \(\mathcal M\) is a well-defined structure of which every
\(\mathcal M_i\) is a substructure. Also \(\prec_A\) is transitive. We show by
induction on \(\varphi(\bar v)\in A\) that for every \(i\in I\) and every
tuple \(\bar a\) from \(M_i\), \(\mathcal M_i\models\varphi(\bar a)\) iff
\(\mathcal M\models\varphi(\bar a)\); every subformula of \(\varphi\) is in
\(A\) and has finitely many free variables. Atomic formulas: \(\mathcal M_i\)
is a substructure of \(\mathcal M\). \(\neg\), \(\bigwedge\), \(\bigvee\): the
induction hypothesis for the constituents. Let \(\varphi=\exists w\,\theta(\bar v,w)\).
A witness \(c\in M_i\) for \(\mathcal M_i\) is a witness for \(\mathcal M\)
by the hypothesis for \(\theta\). Conversely let
\(\mathcal M\models\theta(\bar a,c)\) with \(c\in M_j\), and let \(k\) be
the larger of \(i,j\). The hypothesis for \(\theta\) at \(k\) gives
\(\mathcal M_k\models\theta(\bar a,c)\), so
\(\mathcal M_k\models\varphi(\bar a)\), and \(\mathcal M_i\models\varphi(\bar a)\)
because \(\mathcal M_i=\mathcal M_k\) or \(\mathcal M_i\prec_A\mathcal M_k\)
and \(\varphi\in A\). Let \(\varphi=\forall w\,\theta(\bar v,w)\). If
\(\mathcal M\models\varphi(\bar a)\), then \(\mathcal M\models\theta(\bar a,c)\)
for all \(c\in M_i\), and the hypothesis at \(i\) gives
\(\mathcal M_i\models\varphi(\bar a)\). If \(\mathcal M_i\models\varphi(\bar a)\)
and \(c\in M_j\), let \(k\) be the larger of \(i,j\); then
\(\mathcal M_k\models\varphi(\bar a)\), so \(\mathcal M_k\models\theta(\bar a,c)\),
and the hypothesis at \(k\) gives \(\mathcal M\models\theta(\bar a,c)\).
\end{proof}
\end{annexitem}

\begin{annexitem}{annex:p03-w03-p03-039-uncountable-models}{thm:p03-w03-p03-039-uncountable-models}{When a complete sentence has an uncountable model}
\begin{proof}
If \(M\) is finite, \(\chi_{\mathcal M}\) of
\cref{thm:p03-w03-p03-039-finite-structures} holds in \(\mathcal M\), so
\(\sigma\models\chi_{\mathcal M}\) by completeness, every model of
\(\sigma\) is isomorphic to \(\mathcal M\), and item~1 fails; items~2
and~3 fail because an injective map of a finite set into itself is
surjective. So assume \(M\) is countably infinite.

2\(\Rightarrow\)3: \(A\) consists of \(\mathcal L_{\omega_1,\omega}\)-formulas
with finitely many free variables.

3\(\Rightarrow\)1. Let \(f:M\to M\) be \(A\)-elementary and not surjective.
Its image is the universe of a substructure \(\mathcal M'\) of \(\mathcal M\),
\(f:\mathcal M\to\mathcal M'\) is an isomorphism, and
\(\mathcal M'\prec_A\mathcal M\): for \(\bar a=f(\bar b)\) and
\(\varphi\in A\), \(\mathcal M'\models\varphi(\bar a)\) iff
\(\mathcal M\models\varphi(\bar b)\)
(\cref{thm:p03-w03-p03-039-iso-invariance}) iff
\(\mathcal M\models\varphi(\bar a)\). By recursion on \(\alpha<\omega_1\)
build countable \(\tau\)-structures \(\mathcal M_\alpha\) and isomorphisms
\(h_\alpha:\mathcal M\to\mathcal M_\alpha\) with
\(\mathcal M_\beta\prec_A\mathcal M_\alpha\) for \(\beta<\alpha\) and
\(M_\alpha\subsetneq M_{\alpha+1}\). Put \(\mathcal M_0=\mathcal M\),
\(h_0=\mathrm{id}\). Given \(\mathcal M_\alpha\), the map
\(h_\alpha\circ f^{-1}:\mathcal M'\to\mathcal M_\alpha\) is an isomorphism;
take a set \(Y\) of new elements, disjoint from \(M_\alpha\), in bijection
with \(M\setminus M'\ne\varnothing\), extend \(h_\alpha\circ f^{-1}\) to a
bijection \(g:M\to M_\alpha\cup Y\), and let \(\mathcal M_{\alpha+1}\) be the
structure on \(M_\alpha\cup Y\) that makes \(g\) an isomorphism, with
\(h_{\alpha+1}=g\). Since \(g\) carries \(\mathcal M'\) onto
\(\mathcal M_\alpha\) and \(\mathcal M'\prec_A\mathcal M\),
\(\mathcal M_\alpha\prec_A\mathcal M_{\alpha+1}\), and
\(M_\alpha\subsetneq M_{\alpha+1}\); with transitivity of \(\prec_A\),
\(\mathcal M_\beta\prec_A\mathcal M_{\alpha+1}\) for all \(\beta\le\alpha\).
At a limit \(\lambda<\omega_1\) let \(\mathcal M_\lambda\) be the union; it
is countable, \(\mathcal M_\beta\prec_A\mathcal M_\lambda\) for
\(\beta<\lambda\) by \cref{thm:p03-w03-p03-039-elementary-chain}, and
\(\mathcal M_\lambda\models\sigma\) because \(\mathcal M_0\models\sigma\) and
\(\sigma\in A\). By \cref{thm:p03-w03-p03-039-complete-categorical} there is an
isomorphism \(h_\lambda:\mathcal M\to\mathcal M_\lambda\). Finally let
\(\mathcal M_{\omega_1}\) be the union of all \(\mathcal M_\alpha\),
\(\alpha<\omega_1\). By \cref{thm:p03-w03-p03-039-elementary-chain},
\(\mathcal M_0\prec_A\mathcal M_{\omega_1}\), so
\(\mathcal M_{\omega_1}\models\sigma\). Choosing
\(x_\alpha\in M_{\alpha+1}\setminus M_\alpha\) gives distinct elements
(for \(\alpha<\beta\), \(x_\alpha\in M_\beta\) and \(x_\beta\notin M_\beta\)),
so \(\mathcal M_{\omega_1}\) is uncountable.

1\(\Rightarrow\)2. Let \(\mathcal N\models\sigma\) be uncountable and let
\(\alpha=\operatorname{sr}(\mathcal M)<\omega_1\)
(\cref{thm:p03-w03-p03-039-scott-rank}). The formulas
\(\Phi^{\mathcal M}_{\bar b,\alpha}\), \(\bar b\in M^{<\omega}\), are
countably many \(\mathcal L_{\omega_1,\omega}\)-formulas with finitely many
free variables (\cref{def:p03-w03-p03-039-scott-formulas}). By
\cref{thm:p03-w03-p03-039-countable-fragment-size} there is a countable
fragment \(A'\) containing \(\sigma\) and all of them, whose formulas have
finitely many free variables. Since \(\mathcal M\models\sigma\wedge\Phi^{\mathcal M}\)
and \(\sigma\) is complete, \(\sigma\models\Phi^{\mathcal M}\). By
\cref{thm:p03-w03-p03-039-downward-ls} there is a countable
\(\mathcal N_0\prec_{A'}\mathcal N\); pick \(n\in N\setminus N_0\) and a
countable \(\mathcal N_1\prec_{A'}\mathcal N\) with
\(N_0\cup\{n\}\subseteq N_1\). Then \(\mathcal N_0\) is a substructure of
\(\mathcal N_1\), \(\mathcal N_0\prec_{A'}\mathcal N_1\), and both are models of
\(\sigma\) of cardinality at most \(\aleph_0\), so there are isomorphisms
\(h_0:\mathcal M\to\mathcal N_0\) and \(h_1:\mathcal N_1\to\mathcal M\).

Let \(\bar a\) be a tuple from \(N_0\) and \(\bar b=h_0^{-1}(\bar a)\).
By reflexivity of \(\sim_\alpha\) and
\cref{thm:p03-w03-p03-039-scott-formula-lemma},
\(\mathcal M\models\Phi^{\mathcal M}_{\bar b,\alpha}(\bar b)\); so
\(\mathcal N_0\models\Phi^{\mathcal M}_{\bar b,\alpha}(\bar a)\), hence
\(\mathcal N_1\models\Phi^{\mathcal M}_{\bar b,\alpha}(\bar a)\), hence
\(\mathcal M\models\Phi^{\mathcal M}_{\bar b,\alpha}(h_1(\bar a))\), i.e.
\((\mathcal M,\bar b)\sim_\alpha(\mathcal M,h_1(\bar a))\). By
\cref{thm:p03-w03-p03-039-scott-rank} the two tuples are
\(\sim_\beta\)-related for every \(\beta\), so by
\cref{thm:p03-w03-p03-039-rank-characterization} they satisfy the same
\(\mathcal L_{\infty,\omega}\)-formulas (each formula has some quantifier
rank). For every \(\mathcal L_{\infty,\omega}\)-formula \(\psi(\bar v)\):
\(\mathcal N_0\models\psi(\bar a)\) iff \(\mathcal M\models\psi(\bar b)\) iff
\(\mathcal M\models\psi(h_1(\bar a))\) iff \(\mathcal N_1\models\psi(\bar a)\).
So \(e=h_1\circ h_0:M\to M\) (with \(h_0\) viewed as a map into \(N_1\))
preserves every \(\mathcal L_{\infty,\omega}\)-formula, in particular every
\(\mathcal L_{\omega_1,\omega}\)-formula. Its image is \(h_1[N_0]\), which
omits \(h_1(n)\); so \(e\) is not surjective.
\end{proof}
\end{annexitem}

\begin{annexitem}{annex:p03-w03-p03-039-complete-atomic}{thm:p03-w03-p03-039-complete-atomic}{Complete sentences and atomic models}
\begin{proof}
Let \(\mathcal M\) be the model of \(\varphi\) of cardinality at most
\(\aleph_0\) (\cref{thm:p03-w03-p03-039-finite-structures}).

\emph{Case 1: \(M\) is infinite.} Let \(\alpha=\operatorname{sr}(\mathcal M)<\omega_1\)
and \(T_0=\{\varphi\}\cup\{\forall\bar v\,(\Phi^{\mathcal M}_{\bar b,\alpha}(\bar v)\vee\neg\Phi^{\mathcal M}_{\bar b,\alpha}(\bar v)):\bar b\in M^{<\omega}\}\),
a countable set of \(\mathcal L_{\omega_1,\omega}(\tau)\)-sentences with the
same models as \(\varphi\). Apply \cref{thm:p03-w03-p03-039-chang} to
\(T_0\), with the construction of
\cref{def:p03-w03-p03-039-chang-construction}; the least countable fragment
\(A\) containing \(T_0\) contains each \(\Phi^{\mathcal M}_{\bar b,\alpha}\)
as a subformula. Let \(\tau^*\), \(T^*_C\) and \(\Gamma\) be the vocabulary,
theory and types obtained. The annex proof of
\cref{thm:p03-w03-p03-039-chang} shows: (a) in every model of \(T^*_C\)
omitting \(\Gamma\), \(R_\psi\) holds of a tuple iff \(\psi\) does, for
\(\psi\in A\); (b) for every \(\mathcal N\models T_0\), the \emph{natural
expansion} \(\mathcal N^*\), interpreting each \(R_\psi\) as the set of
tuples satisfying \(\psi\), models \(T^*_C\) and omits \(\Gamma\). Let
\(T^*\) be the complete first-order theory of \(\mathcal M^*\); it contains
\(T^*_C\).

Claim: if \(\mathcal N\models\varphi\), \(\bar a\) is from \(N\), \(\bar b\)
from \(M\) and \((\mathcal M,\bar b)\equiv_{\infty,\omega}(\mathcal N,\bar a)\),
then \(\bar b\) in \(\mathcal M^*\) and \(\bar a\) in \(\mathcal N^*\) satisfy
the same first-order \(\tau^*\)-formulas. Let \(Q\) be the set of finite
maps \(g\) from \(M\) to \(N\), extending \(\bar b\mapsto\bar a\), such that
\((\mathcal M,\bar c)\equiv_{\infty,\omega}(\mathcal N,g(\bar c))\) for an
enumeration \(\bar c\) of \(\operatorname{dom}(g)\). The map
\(\bar b\mapsto\bar a\) is well defined (agreement on \(v_i=v_j\)) and lies in
\(Q\), and \(Q\) has the back and forth properties by the claim in the proof
of (1)\(\Rightarrow\)(4) of \cref{thm:p03-w03-p03-039-karp}. Each \(g\in Q\)
is a partial \(\tau^*\)-embedding: an atomic \(\tau^*\)-formula is either a
\(\tau\)-formula or has the form \(R_\psi(t_1,\ldots,t_k)\) with
\(\tau\)-terms \(t_j(\bar y)\), and in the natural expansions the latter is
equivalent to the \(\mathcal L_{\infty,\omega}(\tau)\)-formula
\(\exists u_1\cdots\exists u_k(\bigwedge_ju_j=t_j\wedge\psi'(\bar u))\),
where \(\psi'\) renames the free variables of \(\psi\), in their fixed order,
to new variables \(u_1,\ldots,u_k\) occurring nowhere else; so its truth
transfers along \(g\). Hence \(Q\) is a partial isomorphism system between the expansions of
\(\mathcal M^*\) and \(\mathcal N^*\) by constants for \(\bar b\) and \(\bar a\),
and \cref{thm:p03-w03-p03-039-karp} gives the claim.

Let \(\mathcal N\models\varphi\). Since \(\varphi\) is complete and
\(\mathcal M\models\varphi\wedge\Phi^{\mathcal M}\),
\(\mathcal N\models\Phi^{\mathcal M}\), so
\(\mathcal N\equiv_{\infty,\omega}\mathcal M\)
(\cref{thm:p03-w03-p03-039-scott}). The claim with empty tuples gives
\(\mathcal N^*\models T^*\). Let \(\bar a\) be from \(N\). Applying the claim
in the proof of \cref{thm:p03-w03-p03-039-karp} once for each entry, there is
\(\bar b\) from \(M\) with \((\mathcal M,\bar b)\equiv_{\infty,\omega}(\mathcal N,\bar a)\).
Let \(\theta(\bar v)\) be the atomic \(\tau^*\)-formula \(R_\psi(\bar v)\) for
\(\psi=\Phi^{\mathcal M}_{\bar b,\alpha}\). If \(\mathcal M^*\models\theta(\bar c)\),
then \((\mathcal M,\bar b)\sim_\alpha(\mathcal M,\bar c)\)
(\cref{thm:p03-w03-p03-039-scott-formula-lemma}), so the two tuples are
\(\sim_\beta\)-related for all \(\beta\)
(\cref{thm:p03-w03-p03-039-scott-rank}), so
\((\mathcal M,\bar b)\equiv_{\infty,\omega}(\mathcal M,\bar c)\)
(\cref{thm:p03-w03-p03-039-rank-characterization}), and by the claim (with
\(\mathcal N=\mathcal M\)) \(\bar c\) and \(\bar b\) have the same
first-order type in \(\mathcal M^*\). Hence
\(\forall\bar v(\theta\to\chi)\in T^*\) for every \(\chi\) in the type of
\(\bar b\), and \(\theta\) isolates that type. By the claim the type of
\(\bar a\) in \(\mathcal N^*\) is the same, so it is isolated. Thus
\(\mathcal N^*\) is an atomic model of \(T^*\) with \(\tau\)-reduct
\(\mathcal N\).

Conversely let \(\mathcal K\) be an atomic model of \(T^*\). Suppose a tuple
\(\bar c\) of \(\mathcal K\) realizes some \(\gamma\in\Gamma\). Its type is
isolated by a formula \(\theta\) with \(\mathcal K\models\theta(\bar c)\), so
\(\exists\bar x\,\theta\in T^*\) and \(\forall\bar x(\theta\to\chi)\in T^*\)
for each \(\chi\in\gamma\); then any \(\bar d\) with
\(\mathcal M^*\models\theta(\bar d)\) realizes \(\gamma\) in \(\mathcal M^*\),
contradicting (b). So \(\mathcal K\models T^*_C\) omits \(\Gamma\), and its
\(\tau\)-reduct is a model of \(T_0\), hence of \(\varphi\), by
\cref{thm:p03-w03-p03-039-chang}.

\emph{Case 2: \(M\) is finite.} Let \(\tau^*=\tau\) and \(T^*\) the complete
first-order theory of \(\mathcal M\). If \(\mathcal N\models\varphi\), then
\(\varphi\models\chi_{\mathcal M}\) by completeness, so
\(\mathcal N\cong\mathcal M\) (\cref{thm:p03-w03-p03-039-finite-structures}).
\(\mathcal M\) is atomic: for a tuple \(\bar a\), there are finitely many
tuples \(\bar b\) of its length; for each \(\bar b\) of a different
first-order type pick \(\psi_{\bar b}\) true of \(\bar a\) and false of
\(\bar b\), and let \(\theta\) be the conjunction of these finitely many
formulas (or \(\bar v=\bar v\) if there are none). Every tuple satisfying
\(\theta\) in \(\mathcal M\) has the type of \(\bar a\), so
\(\forall\bar v(\theta\to\chi)\in T^*\) for every \(\chi\) in that type.
Hence \(\mathcal N\), isomorphic to \(\mathcal M\), is an atomic model of
\(T^*\). Conversely let \(\mathcal K\models T^*\); \(T^*\) says there are
exactly \(n=\lvert M\rvert\) elements. For a finite
\(\tau_0\subseteq\tau\), the formulas \(x_i=x_j\), \(R(x_{i_1},\ldots)\),
\(F(x_{i_1},\ldots)=x_j\), \(c=x_j\) with symbols from \(\tau_0\) and
variables \(x_1,\ldots,x_n\) are finitely many; the first-order sentence
saying that some \(n\) distinct elements exhaust the universe and satisfy
exactly the ones of these formulas that a fixed enumeration of \(M\)
satisfies is in \(T^*\). So the set \(S_{\tau_0}\) of bijections \(M\to K\)
that are \(\tau_0\)-isomorphisms is nonempty and finite, and it shrinks as
\(\tau_0\) grows. Writing \(\tau\) as the union of finite
\(\tau_0\subseteq\tau_1\subseteq\cdots\), the decreasing nonempty finite sets
\(S_{\tau_k}\) have a common member, which is a \(\tau\)-isomorphism. So
\(\mathcal K\cong\mathcal M\models\varphi\).
\end{proof}
\end{annexitem}

\begin{annexitem}{annex:p03-w03-p03-039-countably-many-types}{thm:p03-w03-p03-039-countably-many-types}{Countably many types and countable \(\equiv_{\infty,\omega}\)-representatives}
\begin{proof}
(1) For a tuple \(\bar a\) from \(M\), applying the claim in the proof of
(1)\(\Rightarrow\)(4) of \cref{thm:p03-w03-p03-039-karp} once for each entry
gives \(\bar b\) from \(N\) with
\((\mathcal M,\bar a)\equiv_{\infty,\omega}(\mathcal N,\bar b)\); so the
\(\mathcal L_{\omega_1,\omega}\)-type of \(\bar a\) is realized in
\(\mathcal N\), which has countably many tuples.

For (2) and (3) let \(R_n\) be the countable set of
\(\mathcal L_{\omega_1,\omega}\)-types of \(n\)-tuples realized in
\(\mathcal M\), in the variables \(v_0,\ldots,v_{n-1}\). For distinct
\(p,q\in R_n\) choose \(\psi_{p,q}\in p\setminus q\); one exists, since if
\(\chi\in q\setminus p\), then \(\neg\chi\in p\setminus q\). Put
\(\theta_p=\bigwedge_{q\in R_n,q\ne p}\psi_{p,q}\), a countable conjunction
with free variables among \(v_0,\ldots,v_{n-1}\). An \(n\)-tuple \(\bar c\)
of \(\mathcal M\) satisfies \(\theta_p\) iff its type is \(p\): its type
\(q\) is in \(R_n\), and if \(q\ne p\) then \(\psi_{p,q}\notin q\).

Claim: if \(\bar a,\bar b\) have the same type and \(c\in M\), some \(d\in M\)
gives \(\bar ac\) and \(\bar bd\) the same type. Let \(q\) be the type of
\(\bar ac\). The \(\mathcal L_{\omega_1,\omega}\)-formula
\(\exists v_n\theta_q\) is true of \(\bar a\), hence of \(\bar b\); a witness
\(d\) gives \(\bar bd\) the type \(q\).

(2) Let \(E\) be the set of pairs of tuples with the same type. By induction
on \(\beta\), \(E\subseteq\Gamma_\beta\): at \(0\) because atomic formulas are
in \(\mathcal L_{\omega_1,\omega}\), at successors by the claim (and its
symmetric form), at limits trivially. Every \(\mathcal L_{\omega_1,\omega}\)-formula
has countable quantifier rank (by induction: a countable supremum of
countable ordinals is countable). Let \(\gamma<\omega_1\) be the supremum of
the quantifier ranks of the countably many \(\psi_{p,q}\). If
\((\bar a,\bar b)\in\Gamma_\gamma\) and the types \(p,q\) of \(\bar a,\bar b\)
differ, then \(\bar a\) satisfies \(\psi_{p,q}\) and \(\bar b\) does not,
contradicting \cref{thm:p03-w03-p03-039-rank-characterization}. So
\(\Gamma_\gamma\subseteq E\subseteq\Gamma_{\gamma+1}\subseteq\Gamma_\gamma\),
and \(\operatorname{sr}(\mathcal M)\le\gamma<\omega_1\)
(\cref{thm:p03-w03-p03-039-scott-rank}).

(3) If \(M\) is finite take \(\mathcal N=\mathcal M\). Otherwise, by
\cref{thm:p03-w03-p03-039-countable-fragment-size} let \(A\) be a countable
fragment containing every \(\theta_p\), with witness the variables
\(v_0,v_1,\ldots\), whose formulas have finitely many free variables; it
contains every \(\exists v_n\theta_q\). By
\cref{thm:p03-w03-p03-039-downward-ls} there is a countable
\(\mathcal N\prec_A\mathcal M\). Let \(P\) be the set of finite maps \(g\)
from \(N\) to \(M\) such that \(\bar a\) and \(g(\bar a)\) have the same type
in \(\mathcal M\) for an enumeration \(\bar a\) of \(\operatorname{dom}(g)\)
(this does not depend on the enumeration). The empty map is in \(P\). Each
\(g\in P\) is a partial embedding of \(\mathcal N\) into \(\mathcal M\),
since \(\mathcal N\) is a substructure of \(\mathcal M\) and the types
contain the atomic formulas. Forth: given \(c\in N\), the claim gives
\(d\in M\). Back: given \(d\in M\), let \(q\) be the type of \(g(\bar a)d\);
then \(\mathcal M\models\exists v_n\theta_q(g(\bar a))\), so
\(\mathcal M\models\exists v_n\theta_q(\bar a)\), so
\(\mathcal N\models\exists v_n\theta_q(\bar a)\) as \(\mathcal N\prec_A\mathcal M\);
a witness \(c\in N\) satisfies \(\mathcal M\models\theta_q(\bar a,c)\), so
\(\bar ac\) has type \(q\). In both cases the extended map is a function, by
agreement on the formulas \(v_i=v_n\). So \(P\) is a partial isomorphism
system and \(\mathcal N\equiv_{\infty,\omega}\mathcal M\)
(\cref{thm:p03-w03-p03-039-karp}).
\end{proof}
\end{annexitem}

\begin{annexitem}{annex:p03-w03-p03-039-club-filter}{thm:p03-w03-p03-039-club-filter}{Clubs are closed under countable and diagonal intersections}
\begin{proof}
(3) A union of sets containing \(s\) contains \(s\), and \(a\cup s\) is
countable.

(1) \(\mathcal P_{\omega_1}(X)\) is a club, since a countable union of
countable sets is countable. Let \(\mathcal A_k\), \(k\in\omega\), be clubs
(a finite family is listed with repetitions). The intersection is closed,
since a chain in it is a chain in each \(\mathcal A_k\). For unboundedness
let \(a\in\mathcal P_{\omega_1}(X)\) and fix a bijection
\(\pi:\omega\to\omega\times\omega\). Put \(a_{-1}=a\) and, given
\(a_{n-1}\), choose \(a_n\in\mathcal A_k\) with \(a_{n-1}\subseteq a_n\),
where \(k\) is the first coordinate of \(\pi(n)\). Let \(b=\bigcup_na_n\).
For each \(k\) the indices \(n\) whose \(\pi(n)\) has first coordinate \(k\)
form an infinite set \(n_0<n_1<\cdots\), and \((a_{n_j})_j\) is a chain in
\(\mathcal A_k\) whose union is \(b\), because it is cofinal in the chain
\((a_n)\). So \(b\in\mathcal A_k\) for every \(k\), and \(a\subseteq b\).
Every club is nonempty (apply unboundedness to \(\varnothing\)), so
\(\varnothing\notin D(X)\); \(D(X)\) is closed upwards by definition and
under countable intersections by what was just shown.

(2) Closed: let \(a_0\subseteq a_1\subseteq\cdots\) be in the diagonal
intersection and \(a=\bigcup_na_n\). For \(j\in a\cap I\) there is \(n_0\)
with \(j\in a_n\) for \(n\ge n_0\); then \(a_n\in\mathcal A_j\) for
\(n\ge n_0\), and \(a\), the union of this tail, is in \(\mathcal A_j\).
Unbounded: let \(a_0=a\). Given \(a_n\), the set \(a_n\cap I\) is countable,
so \(\bigcap_{i\in a_n\cap I}\mathcal A_i\) is a club by (1) (it is
\(\mathcal P_{\omega_1}(X)\) when \(a_n\cap I=\varnothing\)); choose
\(a_{n+1}\supseteq a_n\) in it. Let \(b=\bigcup_na_n\). For \(i\in b\cap I\)
there is \(n_0\) with \(i\in a_n\) for \(n\ge n_0\), so \(a_{n+1}\in\mathcal A_i\)
for \(n\ge n_0\), and \(b\in\mathcal A_i\) since \(\mathcal A_i\) is closed.
So \(b\) is in the diagonal intersection and \(a\subseteq b\).
\end{proof}
\end{annexitem}

\begin{annexitem}{annex:p03-w03-p03-039-kueker-game}{thm:p03-w03-p03-039-kueker-game}{Kueker's game}
\begin{proof}
A strategy for Player~II is a function \(s\) from nonempty finite sequences
of elements of \(X\) to \(X\); II follows \(s\) when
\(d_n=s(c_0,\ldots,c_n)\) for all \(n\).

Suppose \(\mathcal A\) includes a club \(\mathcal C\). Fix a bijection
\(\pi:\omega\to\omega\times\omega\) with \(i\le n\) whenever
\(\pi(n)=(i,j)\) (for instance the inverse of
\((i,j)\mapsto(i+j)(i+j+1)/2+j\)). Player~II plays as follows, using a fixed
choice function for every choice. At move \(n\), having seen
\(c_0,\ldots,c_n\), she chooses \(b_n\in\mathcal C\) containing
\(b_{n-1}\cup\{c_0,\ldots,c_n,d_0,\ldots,d_{n-1}\}\) (with
\(b_{-1}=\varnothing\)) and a surjection \(e_n:\omega\to b_n\) (\(b_n\) is
nonempty), and plays \(d_n=e_i(j)\) where \(\pi(n)=(i,j)\). All of this is
determined by \(c_0,\ldots,c_n\), so it is a strategy. In any play following
it, let \(S=\{c_n,d_n:n\in\omega\}\). Each \(b_i\subseteq S\), since
\(e_i(j)\) is played at the move \(\pi^{-1}(i,j)\); and \(S\subseteq\bigcup_nb_n\),
since \(c_n\in b_n\) and \(d_n\in b_i\). So \(S\) is the union of the chain
\((b_n)\) in \(\mathcal C\), \(S\in\mathcal C\subseteq\mathcal A\), and II wins.

Conversely let \(s\) be a winning strategy for II. Let \(\mathcal C\) be the
set of nonempty \(a\in\mathcal P_{\omega_1}(X)\) such that
\(s(c_0,\ldots,c_n)\in a\) whenever all \(c_i\in a\). It is closed, since a
finite sequence from the union of a chain lies in one member. It is
unbounded: given \(a\), let \(a_0=a\cup\{x\}\) for some \(x\in X\) and
\(a_{k+1}=a_k\cup\{s(\bar c):\bar c\text{ a nonempty finite sequence from }a_k\}\),
countable; \(\bigcup_ka_k\in\mathcal C\). If \(a\in\mathcal C\), let Player~I
play an enumeration \(c_0,c_1,\ldots\) of \(a\) (with repetitions if \(a\)
is finite) against \(s\); every answer lies in \(a\), so the set of the play
is \(a\), and \(a\in\mathcal A\) because \(s\) wins. So
\(\mathcal C\subseteq\mathcal A\) and \(\mathcal A\in D(X)\).
\end{proof}
\end{annexitem}

\begin{annexitem}{annex:p03-w03-p03-039-substructure-ae}{thm:p03-w03-p03-039-substructure-ae}{Approximations are substructures almost everywhere}
\begin{proof}
Let \(\mathcal C\) be the set of \(a\in\mathcal P_{\omega_1}(X)\) such that
\(M\cap a\) is nonempty, contains \(c^{\mathcal M}\) for every constant
symbol \(c\), and contains \(F^{\mathcal M}(\bar e)\) for every function
symbol \(F\) and tuple \(\bar e\) from \(M\cap a\). For such \(a\),
\(M\cap a\) is the universe of a substructure. \(\mathcal C\) is closed,
since each condition concerns finitely many elements. It is unbounded:
given \(a\), let \(a_0=a\cup\{m\}\cup\{c^{\mathcal M}:c\in\tau\}\) for some
\(m\in M\), and \(a_{k+1}=a_k\cup\{F^{\mathcal M}(\bar e):F\in\tau,\ \bar e\text{ from }M\cap a_k\}\),
countable because \(\tau\) and \(a_k\) are; \(\bigcup_ka_k\in\mathcal C\),
and all these sets lie in \(X\) because \(M\subseteq X\).
\end{proof}
\end{annexitem}

\begin{annexitem}{annex:p03-w03-p03-039-approximation-nnf}{thm:p03-w03-p03-039-approximation-nnf}{Approximation commutes with negation normal form}
\begin{proof}
The proof of \cref{thm:p03-w03-p03-039-formal-negation} applies verbatim to
indexed families, so \(\models\neg\psi\leftrightarrow{\sim}\psi\) and
\({\sim}\psi\) has the free variables of \(\psi\); by its definition
\({\sim}\psi\) has the index sets of \(\psi\). By induction on \(\psi\), if
\(\psi\) applies \(\neg\) only to atomic formulas, so does \({\sim}\psi\).
Hence, by induction on \(\varphi\), \(\operatorname{nnf}(\varphi)\) applies
\(\neg\) only to atomic formulas, has the free variables and index sets of
\(\varphi\), and \(\models\varphi\leftrightarrow\operatorname{nnf}(\varphi)\);
the step for \(\neg\varphi\) is
\(\operatorname{nnf}(\neg\varphi)={\sim}\operatorname{nnf}(\varphi)\), equivalent
to \(\neg\operatorname{nnf}(\varphi)\) and so to \(\neg\varphi\).

Next, \(({\sim}\psi)^a={\sim}(\psi^a)\) for every \(\psi\), by induction: for
atomic \(\psi\) both sides are \(\neg\psi\); for \(\psi=\neg\chi\) both are
\(\chi^a\); for \(\psi=\bigwedge_{i\in I}\psi_i\),
\(({\sim}\psi)^a=\bigvee_{i\in I\cap a}({\sim}\psi_i)^a=\bigvee_{i\in I\cap a}{\sim}(\psi_i^a)={\sim}(\psi^a)\),
dually for \(\bigvee\); and the quantifier cases are immediate. Then
\(\operatorname{nnf}(\varphi)^a=\operatorname{nnf}(\varphi^a)\) by induction
on \(\varphi\): the case \(\neg\varphi\) is
\(({\sim}\operatorname{nnf}(\varphi))^a={\sim}(\operatorname{nnf}(\varphi)^a)={\sim}\operatorname{nnf}(\varphi^a)=\operatorname{nnf}((\neg\varphi)^a)\),
and the other cases commute with both operations. Finally
\(\models\varphi^a\leftrightarrow\operatorname{nnf}(\varphi^a)\) by the first
part applied to \(\varphi^a\).
\end{proof}
\end{annexitem}

\begin{annexitem}{annex:p03-w03-p03-039-kueker-truth}{thm:p03-w03-p03-039-kueker-truth}{Truth is approximated almost everywhere}
\begin{proof}
Let \(\mathcal S\in D(X)\) be the set of \(a\) for which \(\mathcal M^a\) is a
substructure (\cref{thm:p03-w03-p03-039-substructure-ae}), and
\(\psi=\operatorname{nnf}(\varphi)\); by
\cref{thm:p03-w03-p03-039-approximation-nnf}, \(X\) approximates \(\psi\) and
\(\models\psi^a\leftrightarrow\varphi^a\) for every \(a\). Subformulas of
\(\psi\) have finitely many free variables. For a subformula \(\chi(\bar v)\)
of \(\psi\) and a tuple \(\bar b\) from \(M\) let \(W(\chi,\bar b)\) be the
set of \(a\in\mathcal S\) with \(\bar b\subseteq a\) and
\(\mathcal M^a\models\chi^a(\bar b)\). We show by induction on \(\chi\): if
\(\mathcal M\models\chi(\bar b)\), then \(W(\chi,\bar b)\in D(X)\). The set
\(\{a:\bar b\subseteq a\}\) is a club by
\cref{thm:p03-w03-p03-039-club-filter}~(3).

Atomic and negated atomic \(\chi\): \(\chi^a=\chi\), and for
\(a\in\mathcal S\) containing \(\bar b\), \(\mathcal M^a\models\chi(\bar b)\)
iff \(\mathcal M\models\chi(\bar b)\), as \(\mathcal M^a\) is a substructure.
\(\chi=\bigvee_{i\in I}\chi_i\): choose \(j\) with
\(\mathcal M\models\chi_j(\bar b)\); for \(a\in W(\chi_j,\bar b)\) with
\(j\in a\), \(\chi_j^a\) is a disjunct of \(\chi^a\), so
\(a\in W(\chi,\bar b)\); as \(j\in I\subseteq X\), the set of such \(a\) is
in \(D(X)\). \(\chi=\bigwedge_{i\in I}\chi_i\): for each \(i\in I\) choose a
club \(\mathcal C_i\subseteq W(\chi_i,\bar b)\). As \(I\subseteq X\), the
diagonal intersection \(\triangle_{i\in I}\mathcal C_i\) is a club
(\cref{thm:p03-w03-p03-039-club-filter}~(2)); if \(a\) lies in it, in
\(\mathcal S\), and contains \(\bar b\), then for every \(i\in I\cap a\),
\(a\in\mathcal C_i\) and \(\mathcal M^a\models\chi_i^a(\bar b)\), so
\(a\in W(\chi,\bar b)\). \(\chi=\exists v\,\theta(\bar v,v)\), with \(v\) not
among \(\bar v\): choose \(c\) with \(\mathcal M\models\theta(\bar b,c)\);
then \(W(\theta,\bar bc)\subseteq W(\chi,\bar b)\), since \(c\in M\cap a\) is
a witness. \(\chi=\forall v\,\theta(\bar v,v)\): for each \(c\in M\),
\(\mathcal M\models\theta(\bar b,c)\); choose a club
\(\mathcal C_c\subseteq W(\theta,\bar bc)\). As \(M\subseteq X\),
\(\triangle_{c\in M}\mathcal C_c\) is a club; if \(a\) lies in it, in
\(\mathcal S\), and contains \(\bar b\), then for every \(c\in M\cap a\),
\(\mathcal M^a\models\theta^a(\bar b,c)\), so
\(\mathcal M^a\models\forall v\,\theta^a(\bar b,v)\), i.e.
\(a\in W(\chi,\bar b)\).

If \(\mathcal M\models\varphi\), then \(\mathcal M\models\psi\) and
\(W(\psi,\varnothing)\in D(X)\); on it \(\mathcal M^a\models\varphi^a\). If
\(\mathcal M\not\models\varphi\), the same argument for \(\neg\varphi\)
(which has the index sets of \(\varphi\)) gives
\(\mathcal M^a\models\neg\varphi^a\) a.e.; were also
\(\mathcal M^a\models\varphi^a\) a.e., the two sets would meet, since
\(D(X)\) is a proper filter, which is impossible.
\end{proof}
\end{annexitem}

\begin{annexitem}{annex:p03-w03-p03-039-kueker-corollary}{thm:p03-w03-p03-039-kueker-corollary}{Two special cases}
\begin{proof}
If a set in \(D(X)\) and a club \(\mathcal C\) are given, and a set
\(\mathcal B\) contains their intersection, then \(\mathcal B\in D(X)\)
(\cref{thm:p03-w03-p03-039-club-filter}~(1)). (1) A formula of
\(\mathcal L_{\omega_1,\omega}\) has countably many subformulas (by
induction), so the union \(s\) of the index sets of its subformulas is a
countable subset of \(X\). For \(a\supseteq s\), \(I\cap a=I\) for every such
index set, and \(\varphi^a=\varphi\) by induction. Apply
\cref{thm:p03-w03-p03-039-kueker-truth} with the club \(\{a:s\subseteq a\}\).
(2) If \(M\) is countable, \(\{a:M\subseteq a\}\) is a club, and for its
members \(\mathcal M^a=\mathcal M\). Apply
\cref{thm:p03-w03-p03-039-kueker-truth} with this club.
\end{proof}
\end{annexitem}

\begin{annexitem}{annex:p03-w03-p03-039-kueker}{thm:p03-w03-p03-039-kueker}{Kueker's characterization of \(\equiv_{\infty,\omega}\)}
\begin{proof}
Let \(\mathcal S\) be the set of \(a\) for which \(\mathcal M^a\) and
\(\mathcal N^a\) are both substructures; \(\mathcal S\in D(X)\) by
\cref{thm:p03-w03-p03-039-substructure-ae}.

Forward direction of (1). By \cref{thm:p03-w03-p03-039-karp} there is a
partial isomorphism system \(P\) between \(\mathcal M\) and \(\mathcal N\)
whose members have finite domain. Choose, for \(f\in P\) and \(c\in M\), a
member \(F(f,c)\supseteq f\) of \(P\) with \(c\) in its domain, and for
\(d\in N\) a member \(G(f,d)\supseteq f\) with \(d\) in its image; fix
\(f_0\in P\). Let \(\mathcal C\) be the set of
\(a\supseteq\operatorname{dom}(f_0)\cup\operatorname{img}(f_0)\) such that
for all \(f\in P\) with domain and image inside \(a\), all \(c\in M\cap a\)
and all \(d\in N\cap a\), the domains and images of \(F(f,c)\) and
\(G(f,d)\) lie inside \(a\). \(\mathcal C\) is closed, since each condition
concerns finitely many elements. It is unbounded: from \(a\) build
\(a_0=a\cup\operatorname{dom}(f_0)\cup\operatorname{img}(f_0)\) and let
\(a_{k+1}\) add to \(a_k\) the domains and images of all \(F(f,c)\),
\(G(f,d)\) with \(f\), \(c\), \(d\) as above for \(a_k\); there are
countably many finite maps between subsets of \(a_k\), so \(a_{k+1}\) is
countable, and \(\bigcup_ka_k\in\mathcal C\). For
\(a\in\mathcal C\cap\mathcal S\), the set
\(P_a=\{f\in P:\operatorname{dom}(f)\cup\operatorname{img}(f)\subseteq a\}\)
contains \(f_0\), consists of partial embeddings between the substructures
\(\mathcal M^a\) and \(\mathcal N^a\), and is closed under the extensions
\(F\), \(G\) for \(c\in M\cap a\), \(d\in N\cap a\). So it is a partial
isomorphism system between the countable structures \(\mathcal M^a\) and
\(\mathcal N^a\), which are therefore isomorphic
(\cref{thm:p03-w03-p03-039-countable-back-forth}).

Forward direction of (2). By \cref{thm:p03-w03-p03-039-rank-characterization}
there is a least \(\gamma\) with \(\mathcal M\not\sim_\gamma\mathcal N\).
For tuples \(\bar a\) from \(M\) and \(\bar b\) from \(N\) of equal length
with \((\mathcal M,\bar a)\not\sim_\delta(\mathcal N,\bar b)\) for some
\(\delta\), let \(\delta(\bar a,\bar b)\) be the least such \(\delta\); it is
\(0\) or a successor \(\beta+1\), since a failure at a limit is a failure
below it. In the successor case choose \(t(\bar a,\bar b)\): either some
\(c\in M\) such that \((\mathcal M,\bar ac)\not\sim_\beta(\mathcal N,\bar bd)\)
for all \(d\in N\), or some \(d\in N\) such that
\((\mathcal M,\bar ac)\not\sim_\beta(\mathcal N,\bar bd)\) for all \(c\in M\).
Let \(\mathcal C\) be the set of \(a\) with \(t(\bar a,\bar b)\in a\)
whenever \(t(\bar a,\bar b)\) is defined and \(\bar a,\bar b\) lie in
\(a\); it is a club, by the argument used for \(\mathcal C\) above. Let
\(a\in\mathcal C\cap\mathcal S\) and suppose \(h:\mathcal M^a\to\mathcal N^a\)
is an isomorphism. Starting from the empty tuples, as long as
\(\delta(\bar a,h(\bar a))\) is a successor \(\beta+1\), the element
\(t(\bar a,h(\bar a))\) lies in \(a\); if it is \(c\in M\) extend
\(\bar a\) by \(c\), and if it is \(d\in N\) extend \(\bar a\) by
\(h^{-1}(d)\). In both cases the new pair is \(\bar ac\), \(h(\bar ac)\),
and its value of \(\delta\) is at most \(\beta\). Ordinals do not descend
forever, so after finitely many steps \(\delta(\bar a,h(\bar a))=0\): the
tuples \(\bar a\) and \(h(\bar a)\) differ on an atomic formula in
\(\mathcal M\) and \(\mathcal N\). But \(h\) is an isomorphism between
substructures, so they agree on atomic formulas. Hence no such \(h\) exists,
and \(\mathcal M^a\not\cong\mathcal N^a\) on \(\mathcal C\cap\mathcal S\).

Backward directions. If \(\mathcal M^a\cong\mathcal N^a\) a.e. and
\(\mathcal M\not\equiv_{\infty,\omega}\mathcal N\), the forward direction of
(2) gives \(\mathcal M^a\not\cong\mathcal N^a\) a.e., and the two sets in
\(D(X)\) would be disjoint, which is impossible for a proper filter.
Likewise for (2) using the forward direction of (1).
\end{proof}
\end{annexitem}

\begin{annexitem}{annex:p03-w03-p03-039-kueker-countable}{thm:p03-w03-p03-039-kueker-countable}{Countable structures}
\begin{proof}
The set \(\{a:M\subseteq a\}\) is a club, and on it \(\mathcal M^a=\mathcal M\).
Combine with \cref{thm:p03-w03-p03-039-kueker}~(1) and
\cref{thm:p03-w03-p03-039-club-filter}~(1).
\end{proof}
\end{annexitem}

\begin{annexitem}{annex:p03-w03-p03-039-kueker-free}{thm:p03-w03-p03-039-kueker-free}{Free abelian groups through countable approximations}
\begin{proof}
Take \(X=H\cup F\). Substructures in the vocabulary \(\{+,-,0\}\) are
subgroups. A set \(S\) in an abelian group is independent if no nontrivial
integer combination of finitely many of its members is \(0\).

Suppose \(H\equiv_{\infty,\omega}F\). By
\cref{thm:p03-w03-p03-039-kueker-countable}, \(F\cong H^a\) a.e. Let \(K\)
be a countable subgroup of \(H\); the club \(\{a:K\subseteq a\}\) meets
that set, so \(K\) is a subgroup of a group isomorphic to \(F\), and \(K\)
is free by \cref{thm:p03-w03-p03-039-subgroups-free}. So \(H\) is
\(\aleph_1\)-free. If \(H\) were free on \(n\) generators, it would be
countable, and \(H\cong F\) by \cref{thm:p03-w03-p03-039-karp} and
\cref{thm:p03-w03-p03-039-countable-back-forth}; but \(H/2H\) has \(2^n\)
elements and \(F/2F\) is infinite.

Conversely let \(H\) be \(\aleph_1\)-free and not free on finitely many
generators. Each cyclic subgroup is free, so \(H\) is torsion-free. By
Zorn's lemma there is a maximal independent \(S\subseteq H\). Suppose \(S\)
were finite, \(\lvert S\rvert=n\). For each \(h\in H\) there are \(m\ne0\)
and integers \(k_s\) with \(mh=\sum_sk_ss\) (by maximality, if \(h\notin S\);
trivially if \(h\in S\)); \(\iota(h)=(k_s/m)_{s\in S}\in\mathbb Q^S\) does
not depend on the choice, by independence of \(S\), and \(\iota\) is an
additive map, injective because \(H\) is torsion-free. So \(H\) is
countable, hence free (being \(\aleph_1\)-free), with a basis \(B\) whose
image under \(\iota\) is linearly independent over \(\mathbb Q\) (clear
denominators); so \(\lvert B\rvert\le n\), against the hypothesis. Hence
\(S\) is infinite; let \(S_0\subseteq S\) be countably infinite. For \(a\)
in the club \(\{a:S_0\subseteq a\}\), \(H^a\) is a countable subgroup, so it
is free with some basis \(B_a\); \(B_a\) is infinite, since otherwise the
same map into \(\mathbb Q^{B_a}\) would carry the infinite independent set
\(S_0\) to a linearly independent subset of a finite-dimensional space. So
\(H^a\) is free of countably infinite rank and \(H^a\cong F\). By
\cref{thm:p03-w03-p03-039-kueker-countable}, \(H\equiv_{\infty,\omega}F\).
\end{proof}
\end{annexitem}

\section{Proofs: \texorpdfstring{\(\mathscr L_{\omega_1\omega}\)}{L omega1 omega} and admissible fragments after Nadel}
\label{sec:annex-p03-w03-p03-039-nadel}

Structures are nonempty, except that a substructure generated by the empty
set may be empty, as in \cref{def:p03-w03-p03-039-nadel-nonempty}.

\begin{annexitem}{annex:p03-w03-p03-039-nadel-free-variables}{thm:p03-w03-p03-039-nadel-free-variables}{Subformulas keep finitely many free variables}
\begin{proof}
The immediate constituents of \(\neg\theta\) are \(\theta\); of
\(\bigwedge\Phi\) the members of \(\Phi\); of \(\exists v\theta\) and
\(\forall v\theta\) the formula \(\theta\). The subformulas of \(\varphi\)
form the least set containing \(\varphi\) and closed under immediate
constituents. The set of formulas reachable from \(\varphi\) by a finite
chain \(\varphi=\varphi_0,\varphi_1,\dots,\varphi_k\), each \(\varphi_{i+1}\)
an immediate constituent of \(\varphi_i\), contains \(\varphi\) and is closed
under immediate constituents, so every subformula is reached by such a
chain. Along a chain, the free variables of \(\varphi_{i+1}\) are among those
of \(\varphi_i\), together with \(v\) when \(\varphi_i\) is
\(\exists v\varphi_{i+1}\) or \(\forall v\varphi_{i+1}\): the free variables
of \(\neg\theta\) are those of \(\theta\), those of \(\bigwedge\Phi\) are the
union of those of the members, and those of \(\exists v\theta\) are those of
\(\theta\) other than \(v\). By induction on \(i\), the free variables of
\(\varphi_k\) lie in the free variables of \(\varphi\) together with at most
\(k\) further variables, a finite set.
\end{proof}
\end{annexitem}

\begin{annexitem}{annex:p03-w03-p03-039-nadel-expressive}{thm:p03-w03-p03-039-nadel-expressive}{Five properties expressible in \(\mathscr L_{\omega_1\omega}\)}
\begin{proof}
Each sentence below is built with countable conjunctions and disjunctions
and finitely many quantifiers in each branch, so it lies in
\(\mathscr L_{\omega_1\omega}\).

(1) Let \(\sigma_{\mathrm{fin}}\) be
\(\bigvee_{1\le n<\omega}\exists x_1\cdots\exists x_n\,\forall y\bigvee_{i=1}^{n}y=x_i\).
A structure satisfies the \(n\)-th disjunct exactly when its universe has at
most \(n\) elements; so it satisfies \(\sigma_{\mathrm{fin}}\) exactly when
its universe is finite (it is nonempty).

(2) For \(n\ge1\) let \(x^n\) be the term \(x\cdot x\cdots x\) with \(n\)
factors, and let \(\sigma_{\mathrm{tor}}\) be
\(\forall x\bigvee_{1\le n<\omega}x^n=1\). A group satisfies it exactly when
every element has finite order.

(3) For \(n\ge1\) let \(T_n\) be the countable set of terms in the variables
\(x_1,\dots,x_n\), and let \(\sigma_{\mathrm{fg}}\) be
\(\bigvee_{1\le n<\omega}\exists x_1\cdots\exists x_n\,\forall y\bigvee_{t\in T_n}y=t(x_1,\dots,x_n)\).
In a group \(G\) and for \(a_1,\dots,a_n\in G\), the values
\(t^G(a_1,\dots,a_n)\), \(t\in T_n\), form a set containing the \(a_i\) and
\(1\) and closed under \(\cdot\) and \({}^{-1}\), hence a subgroup; it is
contained in every subgroup containing the \(a_i\), by induction on \(t\).
So it is the subgroup generated by \(a_1,\dots,a_n\), and \(G\) satisfies
\(\sigma_{\mathrm{fg}}\) exactly when some finite nonempty tuple generates
\(G\), that is, when \(G\) is finitely generated (a group is nonempty, so a
generating set can be taken nonempty).

(4) Let \(\sigma_\omega\) be
\(\forall x\bigl(\mathsf Nx\leftrightarrow\bigvee_{n\in\omega}x=\underline n\bigr)\wedge\bigwedge_{m<n<\omega}\neg\,\underline m=\underline n\).
The first conjunct says that \(\mathsf N\) is interpreted as the set of
values of the constants and the second that distinct constants have
distinct values; these are the two clauses defining an \(\omega\)-model.

(5) Let \(\sigma_\Phi\) be
\(\forall x_1\cdots\forall x_k\bigvee_{\varphi\in\Phi}\neg\varphi(x_1,\dots,x_k)\).
It holds exactly when every tuple fails some member of \(\Phi\), that is, when
no tuple satisfies all members of \(\Phi\).
\end{proof}
\end{annexitem}

\begin{annexitem}{annex:p03-w03-p03-039-nadel-well-orders}{thm:p03-w03-p03-039-nadel-well-orders}{Well-orderings in \(\mathscr L_{\omega_1\omega_1}\)}
\begin{proof}
The quantifier \(\exists\{v_n:n\in\omega\}\) holds when some assignment of
elements \(a_n\) to the variables \(v_n\) satisfies the matrix. Let
\((L,<)\) be a linear order.

If \((L,<)\) fails the second conjunct, there are \(a_n\in L\) with
\(a_{n+1}<a_n\) for all \(n\). The nonempty set \(\{a_n:n\in\omega\}\) has no
least element, since below each \(a_n\) lies \(a_{n+1}\); so \((L,<)\) is not
a well-ordering. No choice is used here.

If \((L,<)\) is not a well-ordering, let \(X\subseteq L\) be nonempty without
a least element. For each \(x\in X\) there is \(y\in X\) with \(y<x\), since
\(x\) is not least and the order is linear. By the axiom of dependent choices
(a consequence of the axiom of choice) there is a sequence \(a_0\in X\),
\(a_{n+1}\in X\), \(a_{n+1}<a_n\); assigning \(a_n\) to \(v_n\) satisfies
\(\bigwedge_nv_{n+1}<v_n\), and \((L,<)\) fails \(\sigma_{\mathrm{wo}}\).
\end{proof}
\end{annexitem}

\begin{annexitem}{annex:p03-w03-p03-039-nadel-ordinals}{thm:p03-w03-p03-039-nadel-ordinals}{Each ordinal is characterized by a sentence of \(\mathscr L_{\infty\omega}\)}
\begin{proof}
The definition is a recursion on \(\beta\) in which every conjunction and
disjunction is indexed by the set \(\beta\), and \(\varphi_\beta\) has the one
free variable \(x\); so each \(\varphi_\beta\) and each \(\sigma_\alpha\) is in
\(\mathscr L_{\infty\omega}\). Write \(L_{<a}=\{y\in L:y<a\}\). Distinct
ordinals are not order-isomorphic, so an ordinal isomorphic to a given
well-ordering is unique.

\emph{Claim.} For every \(\beta\), every linear order \((L,<)\) and every
\(a\in L\): \((L,<)\models\varphi_\beta(a)\) if and only if \(L_{<a}\) is
well-ordered of type \(\beta\). The proof is by induction on \(\beta\).

Suppose \(h:(\beta,<)\to L_{<a}\) is an isomorphism. For \(y<a\), \(y=h(\gamma)\)
with \(\gamma<\beta\), and \(L_{<y}=h[\gamma]\) because \(h\) is an order
isomorphism and every element below \(y\) is below \(a\); so \(L_{<y}\) has
type \(\gamma\) and, by the induction hypothesis, \(\varphi_\gamma(y)\)
holds. This gives the first conjunct. For \(\gamma<\beta\), the element
\(y=h(\gamma)\) satisfies \(y<a\) and \(\varphi_\gamma(y)\) by the same
argument, which gives the second conjunct.

Conversely suppose \(\varphi_\beta(a)\) holds. For \(y<a\) the first conjunct
gives \(\gamma<\beta\) with \(\varphi_\gamma(y)\), so by the induction
hypothesis \(L_{<y}\) is well-ordered of a type \(\gamma_y<\beta\), unique by
the remark above. If \(y<y'<a\), let \(k:L_{<y'}\to\gamma_{y'}\) be an
isomorphism; it maps the initial segment \(L_{<y}\) onto the initial segment
of \(\gamma_{y'}\) below \(k(y)\), which is the ordinal \(k(y)\), so
\(\gamma_y=k(y)<\gamma_{y'}\). Thus \(y\mapsto\gamma_y\) is strictly
increasing on the linear order \(L_{<a}\), hence an order embedding into
\(\beta\). By the second conjunct each \(\gamma<\beta\) equals \(\gamma_y\)
for some \(y<a\), so the embedding is onto \(\beta\), and \(L_{<a}\) is
well-ordered of type \(\beta\).

\emph{The sentences.} If \(\mathfrak M\cong(\alpha,<)\), we may take
\(\mathfrak M=(\alpha,<)\) by \cref{thm:p03-w03-p03-039-iso-invariance}. It is a
linear order; each \(\delta<\alpha\) has initial segment \(\delta\), so
\(\varphi_\delta(\delta)\) holds by the claim, giving
\(\forall x\bigvee_{\beta<\alpha}\varphi_\beta(x)\); and for \(\beta<\alpha\)
the element \(\beta\) witnesses \(\exists x\,\varphi_\beta(x)\). Conversely
let \(\mathfrak M\models\sigma_\alpha\). Then \(\mathfrak M\) is a linear
order, and by the claim each \(x\in M\) has \(M_{<x}\) well-ordered of a type
\(\beta_x<\alpha\). The argument of the converse direction of the claim,
with \(M\) in place of \(L_{<a}\) and \(\alpha\) in place of \(\beta\), shows
that \(x\mapsto\beta_x\) is an order isomorphism of \(\mathfrak M\) onto
\((\alpha,<)\), the third conjunct of \(\sigma_\alpha\) giving surjectivity.
\end{proof}
\end{annexitem}

\begin{annexitem}{annex:p03-w03-p03-039-nadel-bridge}{thm:p03-w03-p03-039-nadel-bridge}{Nadel's and Marker's partial isomorphisms agree}
\begin{proof}
(1) For \(D\subseteq M\) let \(\langle D\rangle^{\mathfrak M}\) be the set of
values \(t^{\mathfrak M}[\bar a]\) of terms \(t(\bar x)\) at tuples \(\bar a\)
from \(D\). It contains \(D\) and the values of the constants, is closed
under the functions of \(\mathfrak M\), and lies inside every substructure
containing \(D\) (induction on \(t\)); so it is the universe of the
substructure generated by \(D\).

Suppose \(g\) is an isomorphism of \(\langle D\rangle^{\mathfrak M}\) onto
\(\langle f[D]\rangle^{\mathfrak N}\) extending \(f\), where
\(D=\operatorname{dom}f\). By induction on \(t\),
\(g(t^{\mathfrak M}[\bar a])=t^{\mathfrak N}[f\bar a]\) for tuples \(\bar a\)
from \(D\), because the functions of a substructure are restrictions of those
of the structure and \(g\) commutes with them. An atomic formula is
\(R(t_1,\dots,t_k)\) or \(t_1=t_2\); since \(g\) preserves and reflects the
relations of the generated substructures, which are the restrictions of
\(R^{\mathfrak M}\) and \(R^{\mathfrak N}\), and is injective,
\(\mathfrak M\models\varphi[\bar a]\) if and only if
\(\mathfrak N\models\varphi[f\bar a]\) for every atomic \(\varphi\). So \(f\)
is a partial \(\tau\)-embedding.

Conversely let \(f\) be a partial \(\tau\)-embedding with domain \(D\). Put
\(g(t^{\mathfrak M}[\bar a])=t^{\mathfrak N}[f\bar a]\). For terms \(t,t'\)
and tuples \(\bar a,\bar a'\) from \(D\),
\(t^{\mathfrak M}[\bar a]=t'^{\mathfrak M}[\bar a']\) if and only if
\(\mathfrak M\models(t=t')[\bar a\bar a']\), if and only if
\(\mathfrak N\models(t=t')[f\bar a\,f\bar a']\), if and only if
\(t^{\mathfrak N}[f\bar a]=t'^{\mathfrak N}[f\bar a']\); so \(g\) is well
defined and injective. Its image is \(\langle f[D]\rangle^{\mathfrak N}\). It
commutes with each function symbol \(F\), since
\(F^{\mathfrak M}(t_1^{\mathfrak M}[\bar a],\dots)=F(t_1,\dots)^{\mathfrak M}[\bar a]\),
and with the constants. For a relation symbol \(R\),
\(R^{\mathfrak M}(t_1^{\mathfrak M}[\bar a],\dots)\) holds if and only if
\(\mathfrak M\models R(t_1,\dots)[\bar a]\), if and only if
\(\mathfrak N\models R(t_1,\dots)[f\bar a]\). Taking \(t\) a variable shows
\(g\supseteq f\). So \(g\) is the required isomorphism.

(2) Let \(F\ne\varnothing\) consist of partial isomorphisms, which by (1) are
partial \(\tau\)-embeddings. An \(\omega\)-back and forth set is an
\(n\)-back and forth set, since \(\lvert X\rvert<n\) implies
\(\lvert X\rvert<\omega\). An \(n\)-back and forth set with \(n\ge2\) is a
partial isomorphism system: clauses 1 and 2 of
\cref{def:p03-w03-p03-039-partial-iso} are (i) and (ii) for the one-element
sets \(\{a\}\) and \(\{b\}\). A partial isomorphism system is an
\(\omega\)-back and forth set: given \(f\in F\) and \(X=\{a_1,\dots,a_k\}\),
apply clause 1 \(k\) times to obtain \(f\subseteq g_1\subseteq\cdots\subseteq g_k\)
in \(F\) with \(a_i\in\operatorname{dom}g_i\subseteq\operatorname{dom}g_k\)
(for \(k=0\) take \(g_k=f\)); the case of \(Y\subseteq N\) is the same with
clause 2. The final assertion follows, a partial isomorphism system being
nonempty by definition.
\end{proof}
\end{annexitem}

\begin{annexitem}{annex:p03-w03-p03-039-nadel-karp}{thm:p03-w03-p03-039-nadel-karp}{Karp's theorem (Nadel's Theorem~2.1.1)}
\begin{proof}
First, \(\equiv_{\infty\omega}\) of \cref{def:p03-w03-p03-039-nadel-logics}
and \(\equiv_{\infty,\omega}\) of \cref{def:p03-w03-p03-039-language} are the
same relation. Every sentence in Nadel's sense is a sentence in the sense of
\cref{def:p03-w03-p03-039-language}, with \(\bigvee\Phi\) read as
\(\neg\bigwedge\{\neg\varphi:\varphi\in\Phi\}\). Conversely, replacing
recursively each \(\bigvee_{\varphi\in X}\varphi\) in a formula of
\cref{def:p03-w03-p03-039-language} by
\(\neg\bigwedge\{\neg\varphi^*:\varphi\in X\}\), where \(\varphi^*\) is the
result for \(\varphi\), gives a formula \(\varphi^*\) with the same free
variables and the same satisfaction relation (induction on \(\varphi\)); if
\(\varphi\) is a sentence, \(\varphi^*\) is a sentence all of whose subformulas
have finitely many free variables by
\cref{thm:p03-w03-p03-039-nadel-free-variables}, that is, a sentence in
Nadel's sense. So the two classes of sentences define the same relation of
agreement.

Now \(\mathfrak M\cong_p\mathfrak N\) in the sense of
\cref{def:p03-w03-p03-039-nadel-back-forth}, with a nonempty set, holds if
and only if \(\mathfrak M\cong_p\mathfrak N\) in the sense of
\cref{def:p03-w03-p03-039-partial-iso}, by
\cref{thm:p03-w03-p03-039-nadel-bridge}; and that holds if and only if
\(\mathfrak M\equiv_{\infty,\omega}\mathfrak N\), by the equivalence of
clauses 1 and 2 of \cref{thm:p03-w03-p03-039-karp}.
\end{proof}
\end{annexitem}

\begin{annexitem}{annex:p03-w03-p03-039-nadel-countable}{thm:p03-w03-p03-039-nadel-countable}{The weak form of Scott's theorem (Nadel's Corollary~2.1.2)}
\begin{proof}
By \cref{thm:p03-w03-p03-039-nadel-karp} and
\cref{thm:p03-w03-p03-039-nadel-bridge}, \(\mathfrak M\cong_p\mathfrak N\) in
the sense of \cref{def:p03-w03-p03-039-partial-iso}; since both structures
are countable, \cref{thm:p03-w03-p03-039-countable-back-forth} gives
\(\mathfrak M\cong\mathfrak N\).
\end{proof}
\end{annexitem}

\begin{annexitem}{annex:p03-w03-p03-039-nadel-field-dependencies}{thm:p03-w03-p03-039-nadel-field-dependencies}{The algebra needed for the field back-and-forth argument}
\begin{proof}
\emph{Finite algebraic extensions.} If \(a\) is algebraic over a field
\(K\), its monic polynomial of least positive degree is irreducible:
a proper factorization would give a lower-degree polynomial vanishing
at \(a\), since a field has no zero divisors. Write its degree as
\(d\). Polynomial division reduces every element of \(K[a]\) to a
linear combination of \(1,a,\ldots,a^{d-1}\). For a nonzero polynomial
value \(g(a)\), the irreducible minimal polynomial and \(g\) are
coprime, so the Euclidean algorithm gives polynomials \(u,v\) with
\(ug+v p=1\); evaluating gives \(u(a)g(a)=1\). Thus \(K[a]=K(a)\)
is a field of finite dimension over \(K\).

Adjoining finitely many elements algebraic over \(K\), one at a time,
therefore gives a finite-dimensional field extension. The products of
spanning sets for successive extensions span the composite extension,
so no dimension formula is needed for this conclusion. Every element
\(c\) in a finite-dimensional extension is algebraic: sufficiently
many powers \(1,c,c^2,\ldots\) are linearly dependent over the base.
It follows in particular that the elements algebraic over a fixed
subfield are closed under addition, subtraction, multiplication and
inversion of nonzero elements.

\emph{(F4), transitivity.} Let \(L/K\) be algebraic and let \(c\)
satisfy a nonzero polynomial over \(L\). Its finitely many coefficients
generate a finite-dimensional extension \(K_0/K\). The same polynomial
makes \(c\) algebraic over \(K_0\); thus \(K_0(c)/K_0\), and hence
\(K_0(c)/K\), is finite-dimensional. The preceding power-dependence
argument makes \(c\) algebraic over \(K\).

\emph{(F1), relative algebraic closure.} The preceding closure argument
shows that \(A=\overline K^E\) is a subfield of \(E\), algebraic over
\(K\). A nonconstant polynomial in \(A[X]\) has a root \(c\in E\)
because \(E\) is algebraically closed. This root is algebraic over
\(A\), hence over \(K\) by (F4), and consequently belongs to \(A\).
Every nonconstant polynomial over \(A\) therefore has a root in \(A\).
Dividing off a linear factor and repeating proves that every such
polynomial splits there; \(A\) is algebraically closed.

\emph{(F2), extending a field isomorphism.} Let \(u:K\to K'\) be
an isomorphism, and let \(A,A'\) be algebraic closures of these fields.
Order by extension all embeddings \(h:D\to A'\), where
\(K\subseteq D\subseteq A\) is a subfield and \(h\) extends \(u\).
This is a nonempty set, containing \(u\). The union of a chain has a
subfield as domain and is an embedding extending every member of the
chain. Zorn's lemma gives a maximal \(h\).

If some \(a\in A\) lies outside \(D\), it is algebraic over \(D\).
Let \(p\in D[X]\) be its minimal polynomial. Apply \(h\) to its
coefficients. The resulting irreducible polynomial over \(h[D]\) has
a root \(a'\in A'\). The isomorphism between the quotient fields
\(D[X]/(p)\) and \(h[D][X]/(h(p))\) extends \(h\) to an embedding
\(D(a)\to A'\) taking \(a\) to \(a'\), contradicting maximality.
Hence \(D=A\). Its image \(h[A]\) is algebraically closed and contains
\(K'\). Each element of \(A'\) is algebraic over \(K'\), hence over
\(h[A]\); a polynomial over an algebraically closed field splits into
linear factors, so any of its roots in an extension already lies in
that field. Thus \(h[A]=A'\), as required.

\emph{(F3), a transcendental adjunction.} Evaluation at a transcendental
element \(a\) is an injective map \(K[X]\to K[a]\). The coefficient
isomorphism induced by \(u\), followed by evaluation at \(b\), gives
an isomorphism \(K[a]\to K'[b]\) taking \(a\) to \(b\). It extends
to fraction fields by sending \(g(a)/h(a)\) to
\(u(g)(b)/u(h)(b)\). Nonzero denominators stay nonzero, and equality
of fractions is preserved by cross multiplication. These fraction
fields are \(K(a)\) and \(K'(b)\).

\emph{Exchange.} Suppose \(a\) is algebraic over \(F(b)\) and
transcendental over \(F\). Clear denominators in an algebraic equation
for \(a\) to obtain a nonzero polynomial \(g(X,Y)\in F[X,Y]\) with
\(g(a,b)=0\). Write this polynomial in powers of \(Y\). At least one
of its coefficient polynomials in \(X\) is nonzero; that coefficient
cannot vanish at \(a\), since \(a\) is transcendental over \(F\).
Thus \(g(a,Y)\) is a nonzero polynomial over \(F(a)\) with root
\(b\), proving that \(b\) is algebraic over \(F(a)\).

\emph{The finite-generator bound.} Suppose \(E\) is algebraic over
\(K(b_1,\ldots,b_r)\) and \(a_1,\ldots,a_{r+1}\in E\) are
algebraically independent over \(K\). We replace generators one at a
time while preserving algebraicity of \(E\) over the generated field.
At stage \(i\leq r\), suppose the base list consists of
\(a_1,\ldots,a_{i-1}\) and \(r-i+1\) remaining generators. Set
\(F=K(a_1,\ldots,a_{i-1})\). Independence makes \(a_i\)
transcendental over \(F\), while algebraicity of \(E\) makes it
algebraic after all remaining generators have been adjoined. In any
ordering of that list, take the first generator \(b\) at which
\(a_i\) becomes algebraic. If \(F_0\) is \(F\) with the preceding
generators adjoined, exchange says that \(b\) is algebraic over
\(F_0(a_i)\). Replace \(b\) by \(a_i\) in the list. The compositum
of the old and new generated fields is algebraic over the new one,
since its only additional generator is \(b\). The field \(E\) is
algebraic over this compositum, which contains the old generated field.
By (F4), \(E\) is still algebraic over the new field. After \(r\) steps it is algebraic
over \(K(a_1,\ldots,a_r)\), contradicting independence of
\(a_{r+1}\). The case \(r=0\) is the definition of algebraicity.
Any independent set of more than \(r\) elements would contain such an
\((r+1)\)-tuple, so the bound holds for all independent sets.

Finally, a relative algebraic closure of
\(K(b_1,\ldots,b_r)\) is algebraic over that field. If the ambient
field were algebraic over that closure, (F4) would make it algebraic
over the finitely generated field itself, contradicting the bound
when arbitrarily large finite independent sets exist. This proves
the last assertion and all the algebraic dependencies used below.
\end{proof}
\end{annexitem}

\begin{annexitem}{annex:p03-w03-p03-039-nadel-dlo}{thm:p03-w03-p03-039-nadel-dlo}{Dense orders without endpoints}
\begin{proof}
Let \(F\) be the set of maps \(f\) with finite domain \(D\subseteq M\) and
values in \(N\) such that \(a<a'\) if and only if \(f(a)<f(a')\) for
\(a,a'\in D\). Such an \(f\) is injective and preserves and reflects \(<\)
and \(=\), the only atomic formulas, so it is a partial \(\tau\)-embedding and,
by \cref{thm:p03-w03-p03-039-nadel-bridge}, a partial isomorphism. The
empty map is in \(F\). Let \(f\in F\) with domain
\(a_1<\dots<a_k\) and \(b_i=f(a_i)\), so \(b_1<\dots<b_k\), and let \(a\in M\)
be outside the domain. If \(k=0\) choose any \(b\in N\). Otherwise \(a<a_1\),
or \(a_i<a<a_{i+1}\) for some \(i\), or \(a_k<a\); choose \(b\in N\) with
\(b<b_1\) (\(N\) has no least element), \(b_i<b<b_{i+1}\) (density), or
\(b_k<b\) (no greatest element) respectively. Then
\(f\cup\{(a,b)\}\in F\). The symmetric argument extends \(f\) to put a given
\(b\in N\) into the range. So \(F\) is a partial isomorphism system, and by
\cref{thm:p03-w03-p03-039-nadel-bridge} it is a nonempty \(\omega\)-back and
forth set.
\end{proof}
\end{annexitem}

\begin{annexitem}{annex:p03-w03-p03-039-nadel-acf-finite}{thm:p03-w03-p03-039-nadel-acf-finite}{Distinct finite transcendence degrees}
\begin{proof}
Suppose \(F\) is a nonempty \(\omega\)-back and forth set from \(\mathfrak M\)
to \(\mathfrak N\). The inverses of its members form a nonempty
\(\omega\)-back and forth set from \(\mathfrak N\) to \(\mathfrak M\), since
the inverse of an isomorphism of generated substructures is one; so we may
assume \(m<n\). Choose \(y_1,\dots,y_n\in N\) algebraically independent over
the prime field of \(\mathfrak N\), take \(f\in F\), and by clause (ii) take
\(h\in F\) with \(f\subseteq h\) and \(y_1,\dots,y_n\in\operatorname{ra}h\);
let \(x_i\in M\) with \(h(x_i)=y_i\). By
\cref{thm:p03-w03-p03-039-nadel-bridge}, \(h\) preserves and reflects atomic
formulas at tuples from its domain. For \(k\ge1\) the formula
\(\underbrace{1+\dots+1}_{k}=0\) is atomic (with dummy free variables), so
\(\mathfrak M\) and \(\mathfrak N\) have the same characteristic and
isomorphic prime fields. Every nonzero polynomial over the prime field
becomes, after multiplication by a nonzero integer in characteristic
\(0\), the image of a polynomial \(p\in\mathbb Z[X_1,\dots,X_n]\) that is
nonzero over the prime field and has the same zeros. For such \(p\), the
atomic formula \(p(x_1,\dots,x_n)=0\) fails at \(\bar y\) in \(\mathfrak N\),
hence fails at \(\bar x\) in \(\mathfrak M\). So \(x_1,\dots,x_n\) are
algebraically independent over the prime field of \(\mathfrak M\), and the
transcendence degree of \(\mathfrak M\) is at least \(n>m\), a
contradiction.
\end{proof}
\end{annexitem}

\begin{annexitem}{annex:p03-w03-p03-039-nadel-acf-infinite}{thm:p03-w03-p03-039-nadel-acf-infinite}{Infinite transcendence degree}
\begin{proof}
The facts (F1)--(F4) and the finite-generator bound are proved in
\cref{thm:p03-w03-p03-039-nadel-field-dependencies}.
Let \(P_M,P_N\) be the prime fields; equal characteristic makes them
isomorphic. For a finite tuple \(\bar a\) from \(M\) write
\(K_{\bar a}=\overline{P_M(\bar a)}^{\mathfrak M}\), and similarly
\(K'_{\bar b}\) in \(\mathfrak N\). Let \(F\) be the set of field isomorphisms
\(f:K_{\bar a}\to K'_{\bar b}\) for finite tuples \(\bar a,\bar b\). Each
\(K_{\bar a}\) is a subfield, so it is closed under \(+,-,\cdot\) and contains
\(0,1\); the substructure it generates is itself, and every \(f\in F\) is a
partial isomorphism. By (F1) \(K_\varnothing\) and \(K'_\varnothing\) are
algebraic closures of \(P_M\) and \(P_N\), so by (F2) the isomorphism
\(P_M\cong P_N\) extends to some \(f_0\in F\); thus \(F\ne\varnothing\).

Let \(f:K_{\bar a}\to K'_{\bar b}\) be in \(F\) and \(a\in M\). If
\(a\in K_{\bar a}\), \(f\) itself has \(a\) in its domain. Otherwise \(a\) is
transcendental over \(K_{\bar a}\): an element of \(M\) algebraic over
\(K_{\bar a}\) is algebraic over \(P_M(\bar a)\) by (F4), hence lies in
\(K_{\bar a}\). The field \(K'_{\bar b}\) is algebraic over \(P_N(\bar b)\),
so its transcendence degree is at most the length of \(\bar b\); since
\(\mathfrak N\) has infinite transcendence degree, \(\mathfrak N\) is not
algebraic over \(K'_{\bar b}\) (by (F4) it would otherwise be algebraic over
\(P_N(\bar b)\) and have finite transcendence degree), so some \(b\in N\) is
transcendental over \(K'_{\bar b}\). By (F3) \(f\) extends to
\(f_1:K_{\bar a}(a)\cong K'_{\bar b}(b)\). By (F4), an element is algebraic
over \(K_{\bar a}(a)\) if and only if it is algebraic over \(P_M(\bar a,a)\),
since \(K_{\bar a}(a)\) is algebraic over \(P_M(\bar a,a)\) and contains it;
so \(K_{\bar aa}\) is the set of elements of \(M\) algebraic over
\(K_{\bar a}(a)\), an algebraic closure of it by (F1), and likewise
\(K'_{\bar bb}\) for \(K'_{\bar b}(b)\). By (F2) \(f_1\) extends to an
isomorphism \(g:K_{\bar aa}\to K'_{\bar bb}\), which is in \(F\), extends
\(f\), and has \(a\) in its domain. The symmetric argument puts a given
\(b\in N\) into the range. So \(F\) is a partial isomorphism system; by
\cref{thm:p03-w03-p03-039-nadel-bridge} \(\mathfrak M\cong_p\mathfrak N\),
and by \cref{thm:p03-w03-p03-039-nadel-karp}
\(\mathfrak M\equiv_{\infty\omega}\mathfrak N\).
\end{proof}
\end{annexitem}

\begin{annexitem}{annex:p03-w03-p03-039-nadel-ordinals-distinct}{thm:p03-w03-p03-039-nadel-ordinals-distinct}{Distinct ordinals are distinguished}
\begin{proof}
By \cref{thm:p03-w03-p03-039-nadel-ordinals}, \((\alpha,<)\models\sigma_\alpha\),
and \((\beta,<)\models\sigma_\alpha\) would give \((\beta,<)\cong(\alpha,<)\),
impossible for distinct ordinals. The \(\mathscr L_{\infty\omega}\)-sentence
\(\sigma_\alpha\) separates the two structures.
\end{proof}
\end{annexitem}

\begin{annexitem}{annex:p03-w03-p03-039-nadel-model-existence}{thm:p03-w03-p03-039-nadel-model-existence}{Model existence theorem (Nadel's Theorem~3.1.3)}
\begin{proof}
We prove (ii); (i) is the case \(T=\varnothing\). Since
\(\mathscr L_{\omega\omega}(\tau)\subseteq L_B\) and \(L_B\) is countable,
\(\tau\) is countable, and so are the set of sentences of \(L_B\), the set of
basic terms, the set of formulas \(\varphi(v)\) of \(L_B\) with exactly the
free variable \(v\), and \(T\). Enumerate them, with repetitions allowed, as
\(\theta_0,\theta_1,\dots\); \(t_0,t_1,\dots\); \(\psi_0(v),\psi_1(v),\dots\);
and \(\tau_0,\tau_1,\dots\) (the last list is empty if \(T\) is). Every set
\(\Phi\) with \(\bigwedge\Phi\) or \(\bigvee\Phi\) a sentence of \(L_B\) is a
set of sentences of \(L_B\) (by clause (ii) of the definition of fragment and
the computation of free variables), hence countable; when it is nonempty,
enumerate it as \(\varphi^\Phi_0,\varphi^\Phi_1,\dots\) with repetitions. Fix
a map \(n\mapsto(i_n,j_n,k_n)\) from \(\omega\) onto \(\omega^3\) taking each
value infinitely often.

\emph{Construction.} Let \(s_0=s\). Given \(s_n\in S\), with
\((i,j,k)=(i_n,j_n,k_n)\), perform the following steps in turn, each
producing a member of \(S\) that contains the previous one.
\begin{enumerate}
\item[(a)] If \(T\) is nonempty, add \(\tau_i\) (hypothesis of (ii)).
\item[(b)] If \(\theta_i\) is in the current set: if it is \(\neg\varphi\),
      add \({\sim}\varphi\) (C2); if \(\bigwedge\Phi\) with \(\Phi\ne\varnothing\),
      add \(\varphi^\Phi_j\) (C3); if \(\forall v\varphi(v)\), add
      \(\varphi(c_j)\) (C4); if \(\bigvee\Phi\), add some \(\varphi\in\Phi\)
      as (C5) allows; if \(\exists v\varphi(v)\), add some \(\varphi(c)\) as
      (C6) allows; if \(c=d\) with \(c,d\in C\), add \(d=c\) (C7)(i);
      otherwise do nothing.
\item[(c)] Add some \(e=t_j\) as (C7)(iii) allows.
\item[(d)] If \(\theta_i\) and \(\theta_j\) are in the current set, \(\theta_j\)
      is \(c=t\) with \(c\in C\) and \(t\) basic, and \(\theta_i\) is
      \(\psi_k(t)\), add \(\psi_k(c)\) (C7)(ii).
\end{enumerate}
Call the result \(s_{n+1}\). Let \(s_\omega=\bigcup_ns_n\).

\emph{Closure properties.} If a sentence is in \(s_\omega\), it is in some
\(s_m\) and in every later \(s_n\); since every triple recurs at infinitely
many stages, each step is eventually performed for it. Hence:
(P0) \(T\subseteq s_\omega\);
(P1) no \(\varphi\) and \(\neg\varphi\) both lie in \(s_\omega\), since both would
lie in one \(s_n\in S\), against (C1);
(P2) \(\neg\varphi\in s_\omega\) implies \({\sim}\varphi\in s_\omega\);
(P3) \(\bigwedge\Phi\in s_\omega\) implies \(\Phi\subseteq s_\omega\);
(P4) \(\forall v\varphi(v)\in s_\omega\) implies \(\varphi(c)\in s_\omega\) for
all \(c\in C\);
(P5) \(\bigvee\Phi\in s_\omega\) implies \(\varphi\in s_\omega\) for some
\(\varphi\in\Phi\);
(P6) \(\exists v\varphi(v)\in s_\omega\) implies \(\varphi(c)\in s_\omega\) for
some \(c\in C\);
(P7) \((c=d)\in s_\omega\) implies \((d=c)\in s_\omega\);
(P8) if \(\psi(t)\in s_\omega\) and \((c=t)\in s_\omega\), with \(t\) basic and
\(\psi(v)\) a formula of \(L_B\) with the one free variable \(v\), then
\(\psi(c)\in s_\omega\);
(P9) for each basic \(t\) there is \(e\in C\) with \((e=t)\in s_\omega\).
By (P5), \(\bigvee\varnothing\notin s_\omega\).

\emph{Equality.} For \(c,d\in C\) put \(c\approx d\) if \((c=d)\in s_\omega\).
Reflexivity: by (P9) with \(t=c\) there is \(e\) with \((e=c)\in s_\omega\);
by (P7) \((c=e)\in s_\omega\); by (P8) with \(\psi(v)\) the formula \(c=v\),
\(t=e\), from \(\psi(e)\) and \(c=e\) we get \(\psi(c)\), that is \(c=c\).
Symmetry is (P7). Transitivity: from \(c=d\) and \(d=e\), (P8) with
\(\psi(v)\) the formula \(v=e\) and \(t=d\) gives \(c=e\). Write \([c]\) for
the class of \(c\).

\emph{Replacement.} If \(\psi(\bar d)\in s_\omega\), where \(\psi(v_1,\dots,v_r)\)
is a formula of \(L_B\) and \(c_i\approx d_i\) for \(i\le r\), then
\(\psi(\bar c)\in s_\omega\): by (P8), with \(t=d_i\) and the formula having
\(c_1,\dots,c_{i-1}\) and \(d_{i+1},\dots,d_r\) substituted (in \(L_B\) by
clause (iii) of the definition of fragment), the arguments are replaced one
at a time.

\emph{The canonical structure.} Let \(\mathfrak M\) have universe
\(\{[c]:c\in C\}\). For a relation symbol \(R\), let
\(R^{\mathfrak M}([c_1],\dots,[c_r])\) hold if \(R(c_1,\dots,c_r)\in s_\omega\);
by replacement and symmetry this does not depend on the representatives. For
an \(r\)-ary function symbol \(F\), let
\(F^{\mathfrak M}([c_1],\dots,[c_r])=[e]\) where \((e=F(c_1,\dots,c_r))\in s_\omega\),
which exists by (P9). This is well defined: if also \(c_i\approx d_i\) and
\((e'=F(d_1,\dots,d_r))\in s_\omega\), replacement gives
\((e'=F(c_1,\dots,c_r))\in s_\omega\), and (P8) with \(\psi(v)\) the formula
\(e=v\) and \(t=F(\bar c)\) gives \((e=e')\in s_\omega\). For a constant symbol
\(k\), let \(k^{\mathfrak M}=[e]\) where \((e=k)\in s_\omega\); the same
argument with \(t=k\) shows independence of \(e\), and for \(k=c\in C\),
\(c^{\mathfrak M}=[c]\) by reflexivity. Every element is \([c]=c^{\mathfrak M}\),
so \(\mathfrak M\) is canonical, and it is countable.
\emph{(T)} For every basic \(t\) and \(e\in C\) with \((e=t)\in s_\omega\),
\(t^{\mathfrak M}=[e]\): for \(t=c\in C\), \(e\approx c\) and
\(c^{\mathfrak M}=[c]\); for \(t\) a constant symbol or \(F(c_1,\dots,c_r)\),
the value was defined as \([e']\) for some \(e'\) with \((e'=t)\in s_\omega\),
and the independence just shown gives \([e']=[e]\).

\emph{Atomic sentences.} We show: if \(\alpha\) is an atomic sentence of
\(L_B\) and \(\alpha\in s_\omega\), then \(\mathfrak M\models\alpha\); if
\(\neg\alpha\in s_\omega\), then \(\mathfrak M\models\neg\alpha\). Induction on
the number \(N(\alpha)\) of occurrences in \(\alpha\) of function symbols and of
constant symbols outside \(C\). If \(N(\alpha)=0\), \(\alpha\) is
\(R(c_1,\dots,c_r)\) or \(c=d\) with constants in \(C\); then
\(\alpha\in s_\omega\) gives \(\mathfrak M\models\alpha\) by the definitions,
and \(\neg\alpha\in s_\omega\) gives \(\alpha\notin s_\omega\) by (P1), hence
\(\mathfrak M\not\models\alpha\) by the definitions and the independence of
representatives. If \(N(\alpha)>0\), choose an occurrence of a closed term
not in \(C\) of least height; it is a constant symbol outside \(C\) or of
the form \(F(c_1,\dots,c_r)\) with all \(c_i\in C\), so it is a basic term
\(t\). Let \(\psi(v)\) be \(\alpha\) (respectively \(\neg\alpha\)) with that
occurrence replaced by \(v\), a first-order formula and so in \(L_B\). By
(P9) and (P8) there is \(e\) with \((e=t)\in s_\omega\) and
\(\psi(e)\in s_\omega\), and \(N\) decreases, so by the induction hypothesis
\(\mathfrak M\models\psi(e)\). By (T), \(t^{\mathfrak M}=[e]=e^{\mathfrak M}\),
and replacing a subterm by a term of the same value does not change the
truth value of an atomic formula; so \(\mathfrak M\models\psi(t)\).

\emph{Truth lemma.} Let \(\rho\) be the rank of formulas: \(0\) for atomic
formulas, \(\rho(\theta)+1\) for \(\neg\theta\), \(\exists v\theta\),
\(\forall v\theta\), and \(\sup\{\rho(\theta)+1:\theta\in\Theta\}\) for
\(\bigwedge\Theta\) and \(\bigvee\Theta\). Substituting a constant for a
variable does not change \(\rho\). We prove for every sentence \(\varphi\) of
\(L_B\), by induction on \(\rho(\varphi)\):
(a) \(\varphi\in s_\omega\) implies \(\mathfrak M\models\varphi\);
(c) \({\sim}\varphi\in s_\omega\) implies \(\mathfrak M\models\neg\varphi\);
(b) \(\neg\varphi\in s_\omega\) implies \(\mathfrak M\models\neg\varphi\).
Assume all three for sentences of smaller rank.

(c): if \(\varphi\) is atomic, \({\sim}\varphi=\neg\varphi\) and the atomic case
applies. If \(\varphi=\neg\eta\), \({\sim}\varphi=\eta\), and (a) for \(\eta\)
gives \(\mathfrak M\models\eta\). If \(\varphi=\bigwedge H\),
\({\sim}\varphi=\bigvee\{{\sim}\eta:\eta\in H\}\), so by (P5)
\({\sim}\eta\in s_\omega\) for some \(\eta\in H\), and (c) for \(\eta\) gives
\(\mathfrak M\models\neg\eta\). If \(\varphi=\bigvee H\), by (P3)
\({\sim}\eta\in s_\omega\) for all \(\eta\in H\), and (c) for each \(\eta\)
gives \(\mathfrak M\models\neg\bigvee H\). If \(\varphi=\exists v\eta\),
\({\sim}\varphi=\forall v\neg\eta\); by (P4) \(\neg\eta(c)\in s_\omega\) for
all \(c\in C\), (b) for \(\eta(c)\) gives \(\mathfrak M\models\neg\eta(c)\), and
since \(\mathfrak M\) is canonical, \(\mathfrak M\models\neg\exists v\eta\). If
\(\varphi=\forall v\eta\), \({\sim}\varphi=\exists v\neg\eta\); by (P6)
\(\neg\eta(c)\in s_\omega\) for some \(c\), and (b) gives
\(\mathfrak M\models\neg\forall v\eta\).

(b): by (P2), \(\neg\varphi\in s_\omega\) gives \({\sim}\varphi\in s_\omega\),
and (c) for \(\varphi\) applies.

(a): atomic \(\varphi\) is the atomic case. If \(\varphi=\neg\theta\), (b) for
\(\theta\) applies. If \(\varphi=\bigwedge\Phi\), use (P3) and (a) for the
members; if \(\bigvee\Phi\), (P5) and (a); if \(\exists v\theta\), (P6) and
(a) for \(\theta(c)\); if \(\forall v\theta\), (P4), (a) for every
\(\theta(c)\), and canonicity.

By (a), \(\mathfrak M\) satisfies every member of \(s_\omega\supseteq s\cup T\),
so it is a canonical model of \(T\cup s\), and of \(\bigwedge s\).
\end{proof}
\end{annexitem}

\begin{annexitem}{annex:p03-w03-p03-039-nadel-persistence}{thm:p03-w03-p03-039-nadel-persistence}{\(\Delta_0\)-formulas are absolute and \(\Sigma\)-formulas persistent}
\begin{proof}
Write \(\mathfrak A=\langle A,E,\dots\rangle\) and
\(\mathfrak B=\langle B,F,\dots\rangle\). We prove, by induction on the
formula, that for every tuple \(\bar a\) from \(A\): for \(\Delta_0\)
\(\varphi\), \(\mathfrak A\models\varphi[\bar a]\) if and only if
\(\mathfrak B\models\varphi[\bar a]\); for \(\Sigma\) \(\varphi\),
\(\mathfrak A\models\varphi[\bar a]\) implies \(\mathfrak B\models\varphi[\bar a]\).

\(\Delta_0\): atomic formulas are preserved and reflected because
\(\mathfrak A\) is a substructure of \(\mathfrak B\); the cases \(\neg\),
\(\&\), \(\vee\) follow from the induction hypothesis. For
\(\exists u\in v\,\psi\) with \(v\) assigned \(a\in A\): if
\(\mathfrak A\models\exists u\in v\,\psi\), a witness \(c\in A\) with
\((c,a)\in E\) has \((c,a)\in F\), and \(\mathfrak B\models\psi\) at \(c\) by the
induction hypothesis. If \(\mathfrak B\models\exists u\in v\,\psi\) with
witness \(c\in B\), then \((c,a)\in F\) and \(a\in A\) give \(c\in A\), so
\((c,a)\in E\) since \(E\) is the restriction of \(F\), and the induction
hypothesis gives \(\mathfrak A\models\psi\) at \(c\). The case
\(\forall u\in v\,\psi\) is the same: the \(F\)-predecessors in \(B\) of \(a\)
are exactly its \(E\)-predecessors in \(A\).

\(\Sigma\): \(\Delta_0\)-formulas are absolute, hence persistent; \(\&\) and
\(\vee\) follow from the induction hypothesis. For \(\exists u\in v\,\psi\)
and \(\exists u\,\psi\), a witness in \(A\) remains a witness in \(B\)
(with \((c,a)\in E\subseteq F\) in the bounded case) and the induction
hypothesis applies. For \(\forall u\in v\,\psi\) holding in \(\mathfrak A\)
at \(a\): every \(c\in B\) with \((c,a)\in F\) lies in \(A\) with
\((c,a)\in E\), so \(\mathfrak A\models\psi\) at \(c\), and \(\mathfrak B\models\psi\)
at \(c\) by the induction hypothesis.

Relative to \(T\), with all structures models of \(T\): if \(\varphi\) is
absolute, it is persistent, and if \(\neg\varphi\) holds in \(\mathfrak A\),
then \(\varphi\) fails in \(\mathfrak A\), hence in \(\mathfrak B\), so
\(\neg\varphi\) holds in \(\mathfrak B\). Conversely if \(\varphi\) and
\(\neg\varphi\) are persistent and \(\varphi\) holds in \(\mathfrak B\) but
not in \(\mathfrak A\), then \(\neg\varphi\) holds in \(\mathfrak A\) and so in
\(\mathfrak B\), a contradiction; with persistence of \(\varphi\) this is
absoluteness.
\end{proof}
\end{annexitem}

\begin{annexitem}{annex:p03-w03-p03-039-nadel-hkappa}{thm:p03-w03-p03-039-nadel-hkappa}{\(H(\kappa)\) is admissible for regular \(\kappa\)}
\begin{proof}
Let \(A=H(\kappa)\). For any set \(y\), \(\mathrm{TC}(y)=y\cup\bigcup_{z\in y}\mathrm{TC}(z)\),
since the right side is transitive, contains \(y\) as a subset, and lies in
every transitive set containing \(y\) as a subset. If \(y\subseteq\mathrm{TC}(x)\),
then \(\mathrm{TC}(y)\subseteq\mathrm{TC}(x)\).

\emph{Transitivity.} If \(x\in y\in A\), then \(x\subseteq\mathrm{TC}(y)\), so
\(\lvert\mathrm{TC}(x)\rvert\le\lvert\mathrm{TC}(y)\rvert<\kappa\).

\emph{Extensionality.} If \(x,y\in A\) have the same members in \(A\), they
have the same members, all members lying in \(A\); so \(x=y\).

\emph{Foundation} (the scheme
\(\exists x\varphi(x)\to\exists x[\varphi(x)\,\&\,\forall y\in x\,\neg\varphi(y)]\)).
If \(\langle A,\in\rangle\models\exists x\varphi(x)\) with parameters in
\(A\), the set \(X=\{x\in A:\langle A,\in\rangle\models\varphi(x)\}\) is
nonempty and has an \(\in\)-minimal element \(x_0\) by foundation in the
universe; each \(y\in x_0\) lies in \(A\) and not in \(X\), so \(x_0\) is the
required witness.

\emph{Pairing and union.} \(\mathrm{TC}(\{x,y\})=\{x,y\}\cup\mathrm{TC}(x)\cup\mathrm{TC}(y)\)
has cardinality less than \(\kappa\), and \(\bigcup x\subseteq\mathrm{TC}(x)\)
gives \(\mathrm{TC}(\bigcup x)\subseteq\mathrm{TC}(x)\). So \(\{x,y\}\) and
\(\bigcup x\) lie in \(A\), and since \(A\) is transitive they witness the two
axioms in \(\langle A,\in\rangle\).

\emph{\(\Delta_0\)-Separation.} For \(u\in A\), parameters in \(A\) and a
\(\Delta_0\)-formula \(\varphi\), let
\(v=\{x\in u:\langle A,\in\rangle\models\varphi(x)\}\). Then \(v\subseteq u\),
so \(\mathrm{TC}(v)\subseteq\mathrm{TC}(u)\) and \(v\in A\), and \(v\) witnesses
the axiom.

\emph{\(\Delta_0\)-Collection.} Suppose
\(\langle A,\in\rangle\models\forall x\in u\,\exists y\,\varphi(x,y)\). For each
\(x\in u\) choose, by the axiom of choice, \(y_x\in A\) with
\(\langle A,\in\rangle\models\varphi(x,y_x)\), and let \(v=\{y_x:x\in u\}\).
Then \(\mathrm{TC}(v)=v\cup\bigcup_{x\in u}\mathrm{TC}(y_x)\) is a union of at
most \(\lvert u\rvert+1<\kappa\) sets each of cardinality less than \(\kappa\),
so it has cardinality less than \(\kappa\) because \(\kappa\) is regular (for
\(\kappa=\aleph_0\): a finite union of finite sets). Thus \(v\in A\) and
\(\langle A,\in\rangle\models\forall x\in u\,\exists y\in v\,\varphi(x,y)\).

So \(H(\kappa)\) is admissible. The cardinals \(\aleph_0\) and \(\aleph_1\) are
regular (the latter under the axiom of choice), so \(\mathrm{HF}\) and
\(\mathrm{HC}\) are admissible. Every natural number \(n\) has
\(\mathrm{TC}(n)=n\) finite and lies in \(\mathrm{HF}\), while
\(\mathrm{TC}(\omega)=\omega\) is infinite; since the ordinals of a
transitive set form an initial segment of the ordinals,
\(o(\mathrm{HF})=\omega\).
\end{proof}
\end{annexitem}
